% =====================================================================
%  Anatomía de un LLM — Fase 0 — Módulo 0.9
%  Estadística
%  Compilar: latexmk -r ../latexmkrc P09-estadistica.tex
% =====================================================================
\documentclass[11pt,a4paper,twoside,openany]{book}
\usepackage{../estilo-apuntes}

\newcommand{\ApuntesCurso}{Anatomía de un LLM}
\newcommand{\ApuntesModulo}{Módulo 0.9}

% Notación propia del módulo (compatible con el módulo 0.8)
\newcommand{\PP}{\mathbb{P}}
\newcommand{\N}{\mathbb{N}}
\DeclareMathOperator{\Var}{Var}
\DeclareMathOperator{\Cov}{Cov}
\DeclareMathOperator{\Cor}{Cor}
\DeclareMathOperator{\ECM}{ECM}
\DeclareMathOperator{\sesgo}{sesgo}
\DeclareMathOperator{\ee}{ee}
\DeclareMathOperator{\RIC}{RIC}
\DeclareMathOperator{\Bin}{Bin}
\DeclareMathOperator{\Ber}{Ber}
\DeclareMathOperator{\Pois}{Pois}
\DeclareMathOperator{\Expo}{Exp}
\DeclareMathOperator{\Unif}{U}
\DeclareMathOperator{\Geom}{Geom}
\DeclareMathOperator*{\argmax}{arg\,max}
\DeclareMathOperator*{\argmin}{arg\,min}
\newcommand{\Nor}{\mathcal{N}}
\newcommand{\Lik}{\mathcal{L}}
\newcommand{\Xbar}{\overline{X}}
\newcommand{\xbar}{\bar{x}}
\newcommand{\ybar}{\bar{y}}
\newcommand{\iid}{i.i.d.}
\newcommand{\ind}{\mathbf{1}}
\newcommand{\convp}{\xrightarrow{\;\PP\;}}

% Anclas de hyperref únicas para los enunciados (numerados por tema)
\providecommand{\theHdefinicion}{}\renewcommand{\theHdefinicion}{\thechapter.\arabic{definicion}}
\providecommand{\theHteorema}{}\renewcommand{\theHteorema}{\thechapter.\arabic{definicion}}
\providecommand{\theHproposicion}{}\renewcommand{\theHproposicion}{\thechapter.\arabic{definicion}}
\providecommand{\theHlema}{}\renewcommand{\theHlema}{\thechapter.\arabic{definicion}}
\providecommand{\theHcorolario}{}\renewcommand{\theHcorolario}{\thechapter.\arabic{definicion}}
\providecommand{\theHejemplo}{}\renewcommand{\theHejemplo}{\thechapter.\arabic{definicion}}
\providecommand{\theHobservacion}{}\renewcommand{\theHobservacion}{\thechapter.\arabic{definicion}}

\hypersetup{pdftitle={Módulo 0.9. Estadística (resumen)},pdfauthor={Anatomía de un LLM}}

\makeatletter
\newlabel{ej:t1-tokens}{{1.2}{1}{Datos de referencia}{ejemplo.1.2}{}}
\newlabel{ej:t1-tokens@cref}{{[ejemplo][2][]1.2}{[1][1][]1}}
\newlabel{ej:t1-media-mediana}{{1.4}{2}{}{ejemplo.1.4}{}}
\newlabel{ej:t1-media-mediana@cref}{{[ejemplo][4][]1.4}{[1][2][]2}}
\newlabel{ej:t1-cuartiles}{{1.8}{3}{}{ejemplo.1.8}{}}
\newlabel{ej:t1-cuartiles@cref}{{[ejemplo][8][]1.8}{[1][3][]3}}
\newlabel{fig:t1-hist-caja}{{1.1}{3}{Histograma (clases de amplitud 10) y diagrama de caja de las 31 notas de examen del ejemplo~2.33 de~\cite {openstax}: mínimo 33, $Q_1=61$, mediana 73, $Q_3=90$, máximo 100. La cola izquierda más larga se aprecia en ambos gráficos}{figure.1.1}{}}
\newlabel{fig:t1-hist-caja@cref}{{[figure][1][1]1.1}{[1][3][]3}}
\newlabel{ej:t1-varianza}{{1.13}{4}{}{ejemplo.1.13}{}}
\newlabel{ej:t1-varianza@cref}{{[ejemplo][13][]1.13}{[1][4][]4}}
\newlabel{ej:t1-r}{{1.18}{5}{}{ejemplo.1.18}{}}
\newlabel{ej:t1-r@cref}{{[ejemplo][18][]1.18}{[1][5][]5}}
\newlabel{obs:t1-errores}{{1.19}{5}{Errores frecuentes}{observacion.1.19}{}}
\newlabel{obs:t1-errores@cref}{{[observacion][19][]1.19}{[1][5][]5}}
\newlabel{eje:t1-vf}{{1.1}{6}{}{ejercicio.1.1}{}}
\newlabel{eje:t1-vf@cref}{{[ejercicio][1][1]1.1}{[1][5][]6}}
\newlabel{eje:t1-basico}{{1.2}{6}{}{ejercicio.1.2}{}}
\newlabel{eje:t1-basico@cref}{{[ejercicio][2][1]1.2}{[1][6][]6}}
\newlabel{eje:t1-edades}{{1.3}{6}{}{ejercicio.1.3}{}}
\newlabel{eje:t1-edades@cref}{{[ejercicio][3][1]1.3}{[1][6][]6}}
\newlabel{eje:t1-cajas}{{1.4}{6}{}{ejercicio.1.4}{}}
\newlabel{eje:t1-cajas@cref}{{[ejercicio][4][1]1.4}{[1][6][]6}}
\newlabel{eje:t1-agrupados}{{1.5}{6}{}{ejercicio.1.5}{}}
\newlabel{eje:t1-agrupados@cref}{{[ejercicio][5][1]1.5}{[1][6][]6}}
\newlabel{eje:t1-abreviada}{{1.6}{6}{}{ejercicio.1.6}{}}
\newlabel{eje:t1-abreviada@cref}{{[ejercicio][6][1]1.6}{[1][6][]6}}
\newlabel{eje:t1-celsius}{{1.7}{6}{}{ejercicio.1.7}{}}
\newlabel{eje:t1-celsius@cref}{{[ejercicio][7][1]1.7}{[1][6][]6}}
\newlabel{eje:t1-zscore}{{1.8}{6}{}{ejercicio.1.8}{}}
\newlabel{eje:t1-zscore@cref}{{[ejercicio][8][1]1.8}{[1][6][]6}}
\newlabel{eje:t1-cov-tabla}{{1.9}{6}{}{ejercicio.1.9}{}}
\newlabel{eje:t1-cov-tabla@cref}{{[ejercicio][9][1]1.9}{[1][6][]6}}
\newlabel{eje:t1-agrupar}{{1.10}{7}{}{ejercicio.1.10}{}}
\newlabel{eje:t1-agrupar@cref}{{[ejercicio][10][1]1.10}{[1][7][]7}}
\newlabel{eje:t1-mediana-min}{{1.11}{7}{}{ejercicio.1.11}{}}
\newlabel{eje:t1-mediana-min@cref}{{[ejercicio][11][1]1.11}{[1][7][]7}}
\newlabel{eje:t1-r-invariante}{{1.12}{7}{}{ejercicio.1.12}{}}
\newlabel{eje:t1-r-invariante@cref}{{[ejercicio][12][1]1.12}{[1][7][]7}}
\newlabel{eje:t1-chebyshev}{{1.13}{7}{}{ejercicio.1.13}{}}
\newlabel{eje:t1-chebyshev@cref}{{[ejercicio][13][1]1.13}{[1][7][]7}}
\newlabel{eje:t1-welford}{{1.14}{7}{}{ejercicio.1.14}{}}
\newlabel{eje:t1-welford@cref}{{[ejercicio][14][1]1.14}{[1][7][]7}}
\newlabel{eje:t1-varianza-pares}{{1.15}{7}{}{ejercicio.1.15}{}}
\newlabel{eje:t1-varianza-pares@cref}{{[ejercicio][15][1]1.15}{[1][7][]7}}
\newlabel{ej:t2-proporcion}{{2.7}{10}{Proporción muestral}{ejemplo.2.7}{}}
\newlabel{ej:t2-proporcion@cref}{{[ejemplo][7][]2.7}{[1][10][]10}}
\newlabel{obs:t2-errores}{{2.15}{12}{Errores frecuentes}{observacion.2.15}{}}
\newlabel{obs:t2-errores@cref}{{[observacion][15][]2.15}{[1][11][]12}}
\newlabel{eje:t2-conceptos}{{2.1}{12}{}{ejercicio.2.1}{}}
\newlabel{eje:t2-conceptos@cref}{{[ejercicio][1][2]2.1}{[1][12][]12}}
\newlabel{eje:t2-enumeracion}{{2.2}{12}{}{ejercicio.2.2}{}}
\newlabel{eje:t2-enumeracion@cref}{{[ejercicio][2][2]2.2}{[1][12][]12}}
\newlabel{eje:t2-tres-estimadores}{{2.3}{12}{}{ejercicio.2.3}{}}
\newlabel{eje:t2-tres-estimadores@cref}{{[ejercicio][3][2]2.3}{[1][12][]12}}
\newlabel{eje:t2-poisson}{{2.4}{12}{}{ejercicio.2.4}{}}
\newlabel{eje:t2-poisson@cref}{{[ejercicio][4][2]2.4}{[1][12][]12}}
\newlabel{eje:t2-expo-sesgo}{{2.5}{12}{}{ejercicio.2.5}{}}
\newlabel{eje:t2-expo-sesgo@cref}{{[ejercicio][5][2]2.5}{[1][12][]12}}
\newlabel{eje:t2-sumas-solapadas}{{2.6}{13}{}{ejercicio.2.6}{}}
\newlabel{eje:t2-sumas-solapadas@cref}{{[ejercicio][6][2]2.6}{[1][13][]13}}
\newlabel{eje:t2-benchmark}{{2.7}{13}{}{ejercicio.2.7}{}}
\newlabel{eje:t2-benchmark@cref}{{[ejercicio][7][2]2.7}{[1][13][]13}}
\newlabel{eje:t2-combinar}{{2.8}{13}{}{ejercicio.2.8}{}}
\newlabel{eje:t2-combinar@cref}{{[ejercicio][8][2]2.8}{[1][13][]13}}
\newlabel{eje:t2-chebyshev-n}{{2.9}{13}{}{ejercicio.2.9}{}}
\newlabel{eje:t2-chebyshev-n@cref}{{[ejercicio][9][2]2.9}{[1][13][]13}}
\newlabel{eje:t2-ecm-normal}{{2.10}{13}{}{ejercicio.2.10}{}}
\newlabel{eje:t2-ecm-normal@cref}{{[ejercicio][10][2]2.10}{[1][13][]13}}
\newlabel{eje:t2-uniforme}{{2.11}{13}{}{ejercicio.2.11}{}}
\newlabel{eje:t2-uniforme@cref}{{[ejercicio][11][2]2.11}{[1][13][]13}}
\newlabel{eje:t2-max-consistente}{{2.12}{13}{}{ejercicio.2.12}{}}
\newlabel{eje:t2-max-consistente@cref}{{[ejercicio][12][2]2.12}{[1][13][]13}}
\newlabel{eje:t2-cov-desv}{{2.13}{13}{}{ejercicio.2.13}{}}
\newlabel{eje:t2-cov-desv@cref}{{[ejercicio][13][2]2.13}{[1][13][]13}}
\newlabel{eje:t2-s-sesgo}{{2.14}{13}{}{ejercicio.2.14}{}}
\newlabel{eje:t2-s-sesgo@cref}{{[ejercicio][14][2]2.14}{[1][13][]13}}
\newlabel{eje:t2-passk}{{2.15}{13}{}{ejercicio.2.15}{}}
\newlabel{eje:t2-passk@cref}{{[ejercicio][15][2]2.15}{[1][13][]13}}
\newlabel{ej:t3-tcl}{{3.6}{16}{}{ejemplo.3.6}{}}
\newlabel{ej:t3-tcl@cref}{{[ejemplo][6][]3.6}{[1][16][]16}}
\newlabel{ej:t3-binomial}{{3.9}{17}{}{ejemplo.3.9}{}}
\newlabel{ej:t3-binomial@cref}{{[ejemplo][9][]3.9}{[1][17][]17}}
\newlabel{obs:t3-errores}{{3.13}{18}{Errores frecuentes}{observacion.3.13}{}}
\newlabel{obs:t3-errores@cref}{{[observacion][13][]3.13}{[1][17][]18}}
\newlabel{eje:t3-vf}{{3.1}{18}{}{ejercicio.3.1}{}}
\newlabel{eje:t3-vf@cref}{{[ejercicio][1][3]3.1}{[1][18][]18}}
\newlabel{eje:t3-tabla}{{3.2}{18}{}{ejercicio.3.2}{}}
\newlabel{eje:t3-tabla@cref}{{[ejercicio][2][3]3.2}{[1][18][]18}}
\newlabel{eje:t3-normal-basico}{{3.3}{18}{}{ejercicio.3.3}{}}
\newlabel{eje:t3-normal-basico@cref}{{[ejercicio][3][3]3.3}{[1][18][]18}}
\newlabel{eje:t3-openstax-media}{{3.4}{18}{}{ejercicio.3.4}{}}
\newlabel{eje:t3-openstax-media@cref}{{[ejercicio][4][3]3.4}{[1][18][]18}}
\newlabel{eje:t3-openstax-suma}{{3.5}{18}{}{ejercicio.3.5}{}}
\newlabel{eje:t3-openstax-suma@cref}{{[ejercicio][5][3]3.5}{[1][18][]18}}
\newlabel{eje:t3-elecciones}{{3.6}{19}{}{ejercicio.3.6}{}}
\newlabel{eje:t3-elecciones@cref}{{[ejercicio][6][3]3.6}{[1][19][]19}}
\newlabel{eje:t3-redondeo}{{3.7}{19}{}{ejercicio.3.7}{}}
\newlabel{eje:t3-redondeo@cref}{{[ejercicio][7][3]3.7}{[1][19][]19}}
\newlabel{eje:t3-binomial}{{3.8}{19}{}{ejercicio.3.8}{}}
\newlabel{eje:t3-binomial@cref}{{[ejercicio][8][3]3.8}{[1][19][]19}}
\newlabel{eje:t3-benchmark}{{3.9}{19}{}{ejercicio.3.9}{}}
\newlabel{eje:t3-benchmark@cref}{{[ejercicio][9][3]3.9}{[1][19][]19}}
\newlabel{eje:t3-tamano}{{3.10}{19}{}{ejercicio.3.10}{}}
\newlabel{eje:t3-tamano@cref}{{[ejercicio][10][3]3.10}{[1][19][]19}}
\newlabel{eje:t3-t-cuantiles}{{3.11}{19}{}{ejercicio.3.11}{}}
\newlabel{eje:t3-t-cuantiles@cref}{{[ejercicio][11][3]3.11}{[1][19][]19}}
\newlabel{eje:t3-lineal-normal}{{3.12}{19}{}{ejercicio.3.12}{}}
\newlabel{eje:t3-lineal-normal@cref}{{[ejercicio][12][3]3.12}{[1][19][]19}}
\newlabel{eje:t3-chi1}{{3.13}{19}{}{ejercicio.3.13}{}}
\newlabel{eje:t3-chi1@cref}{{[ejercicio][13][3]3.13}{[1][19][]19}}
\newlabel{eje:t3-chi-momentos}{{3.14}{19}{}{ejercicio.3.14}{}}
\newlabel{eje:t3-chi-momentos@cref}{{[ejercicio][14][3]3.14}{[1][19][]19}}
\newlabel{eje:t3-t-simetrica}{{3.15}{20}{}{ejercicio.3.15}{}}
\newlabel{eje:t3-t-simetrica@cref}{{[ejercicio][15][3]3.15}{[1][19][]20}}
\newlabel{ej:t4-momentos}{{4.2}{21}{Momentos en los modelos básicos}{ejemplo.4.2}{}}
\newlabel{ej:t4-momentos@cref}{{[ejemplo][2][]4.2}{[1][21][]21}}
\newlabel{ej:t4-uniforme}{{4.6}{23}{Un caso sin derivadas}{ejemplo.4.6}{}}
\newlabel{ej:t4-uniforme@cref}{{[ejemplo][6][]4.6}{[1][22][]23}}
\newlabel{ej:t4-t}{{4.13}{24}{}{ejemplo.4.13}{}}
\newlabel{ej:t4-t@cref}{{[ejemplo][13][]4.13}{[1][24][]24}}
\newlabel{ej:t4-benchmark}{{4.16}{25}{Exactitud en un benchmark}{ejemplo.4.16}{}}
\newlabel{ej:t4-benchmark@cref}{{[ejemplo][16][]4.16}{[1][25][]25}}
\newlabel{obs:t4-errores}{{4.17}{25}{Errores frecuentes}{observacion.4.17}{}}
\newlabel{obs:t4-errores@cref}{{[observacion][17][]4.17}{[1][25][]25}}
\newlabel{eje:t4-interpretacion}{{4.1}{26}{}{ejercicio.4.1}{}}
\newlabel{eje:t4-interpretacion@cref}{{[ejercicio][1][4]4.1}{[1][25][]26}}
\newlabel{eje:t4-momentos-basico}{{4.2}{26}{}{ejercicio.4.2}{}}
\newlabel{eje:t4-momentos-basico@cref}{{[ejercicio][2][4]4.2}{[1][26][]26}}
\newlabel{eje:t4-geometrica}{{4.3}{26}{}{ejercicio.4.3}{}}
\newlabel{eje:t4-geometrica@cref}{{[ejercicio][3][4]4.3}{[1][26][]26}}
\newlabel{eje:t4-gamma}{{4.4}{26}{}{ejercicio.4.4}{}}
\newlabel{eje:t4-gamma@cref}{{[ejercicio][4][4]4.4}{[1][26][]26}}
\newlabel{eje:t4-uniforme-ab}{{4.5}{26}{}{ejercicio.4.5}{}}
\newlabel{eje:t4-uniforme-ab@cref}{{[ejercicio][5][4]4.5}{[1][26][]26}}
\newlabel{eje:t4-poisson-datos}{{4.6}{26}{}{ejercicio.4.6}{}}
\newlabel{eje:t4-poisson-datos@cref}{{[ejercicio][6][4]4.6}{[1][26][]26}}
\newlabel{eje:t4-soda}{{4.7}{26}{}{ejercicio.4.7}{}}
\newlabel{eje:t4-soda@cref}{{[ejercicio][7][4]4.7}{[1][26][]26}}
\newlabel{eje:t4-alturas}{{4.8}{26}{}{ejercicio.4.8}{}}
\newlabel{eje:t4-alturas@cref}{{[ejercicio][8][4]4.8}{[1][26][]26}}
\newlabel{eje:t4-openstax-z}{{4.9}{27}{}{ejercicio.4.9}{}}
\newlabel{eje:t4-openstax-z@cref}{{[ejercicio][9][4]4.9}{[1][27][]27}}
\newlabel{eje:t4-proporcion}{{4.10}{27}{}{ejercicio.4.10}{}}
\newlabel{eje:t4-proporcion@cref}{{[ejercicio][10][4]4.10}{[1][27][]27}}
\newlabel{eje:t4-margen}{{4.11}{27}{}{ejercicio.4.11}{}}
\newlabel{eje:t4-margen@cref}{{[ejercicio][11][4]4.11}{[1][27][]27}}
\newlabel{eje:t4-emv-poisson}{{4.12}{27}{}{ejercicio.4.12}{}}
\newlabel{eje:t4-emv-poisson@cref}{{[ejercicio][12][4]4.12}{[1][27][]27}}
\newlabel{eje:t4-emv-expo}{{4.13}{27}{}{ejercicio.4.13}{}}
\newlabel{eje:t4-emv-expo@cref}{{[ejercicio][13][4]4.13}{[1][27][]27}}
\newlabel{eje:t4-laplace}{{4.14}{27}{}{ejercicio.4.14}{}}
\newlabel{eje:t4-laplace@cref}{{[ejercicio][14][4]4.14}{[1][27][]27}}
\newlabel{eje:t4-regla3}{{4.15}{27}{}{ejercicio.4.15}{}}
\newlabel{eje:t4-regla3@cref}{{[ejercicio][15][4]4.15}{[1][27][]27}}
\newlabel{eje:t4-conservador}{{4.16}{27}{}{ejercicio.4.16}{}}
\newlabel{eje:t4-conservador@cref}{{[ejercicio][16][4]4.16}{[1][27][]27}}
\newlabel{ej:t5-z}{{5.6}{30}{Contraste $z$ unilateral}{ejemplo.5.6}{}}
\newlabel{ej:t5-z@cref}{{[ejemplo][6][]5.6}{[1][29][]30}}
\newlabel{ej:t5-dado}{{5.9}{30}{¿Es equilibrado un dado?}{ejemplo.5.9}{}}
\newlabel{ej:t5-dado@cref}{{[ejemplo][9][]5.9}{[1][30][]30}}
\newlabel{obs:t5-errores}{{5.14}{31}{Errores frecuentes}{observacion.5.14}{}}
\newlabel{obs:t5-errores@cref}{{[observacion][14][]5.14}{[1][31][]31}}
\newlabel{eje:t5-vf}{{5.1}{32}{}{ejercicio.5.1}{}}
\newlabel{eje:t5-vf@cref}{{[ejercicio][1][5]5.1}{[1][32][]32}}
\newlabel{eje:t5-moneda}{{5.2}{32}{}{ejercicio.5.2}{}}
\newlabel{eje:t5-moneda@cref}{{[ejercicio][2][5]5.2}{[1][32][]32}}
\newlabel{eje:t5-poligrafo}{{5.3}{32}{}{ejercicio.5.3}{}}
\newlabel{eje:t5-poligrafo@cref}{{[ejercicio][3][5]5.3}{[1][32][]32}}
\newlabel{eje:t5-radar}{{5.4}{32}{}{ejercicio.5.4}{}}
\newlabel{eje:t5-radar@cref}{{[ejercicio][4][5]5.4}{[1][32][]32}}
\newlabel{eje:t5-uniforme}{{5.5}{33}{}{ejercicio.5.5}{}}
\newlabel{eje:t5-uniforme@cref}{{[ejercicio][5][5]5.5}{[1][33][]33}}
\newlabel{eje:t5-t-test}{{5.6}{33}{}{ejercicio.5.6}{}}
\newlabel{eje:t5-t-test@cref}{{[ejercicio][6][5]5.6}{[1][33][]33}}
\newlabel{eje:t5-openstax-z}{{5.7}{33}{}{ejercicio.5.7}{}}
\newlabel{eje:t5-openstax-z@cref}{{[ejercicio][7][5]5.7}{[1][33][]33}}
\newlabel{eje:t5-openstax-prop}{{5.8}{33}{}{ejercicio.5.8}{}}
\newlabel{eje:t5-openstax-prop@cref}{{[ejercicio][8][5]5.8}{[1][33][]33}}
\newlabel{eje:t5-chi-dias}{{5.9}{33}{}{ejercicio.5.9}{}}
\newlabel{eje:t5-chi-dias@cref}{{[ejercicio][9][5]5.9}{[1][33][]33}}
\newlabel{eje:t5-benford}{{5.10}{33}{}{ejercicio.5.10}{}}
\newlabel{eje:t5-benford@cref}{{[ejercicio][10][5]5.10}{[1][33][]33}}
\newlabel{eje:t5-dos-modelos}{{5.11}{33}{}{ejercicio.5.11}{}}
\newlabel{eje:t5-dos-modelos@cref}{{[ejercicio][11][5]5.11}{[1][33][]33}}
\newlabel{eje:t5-multiples}{{5.12}{33}{}{ejercicio.5.12}{}}
\newlabel{eje:t5-multiples@cref}{{[ejercicio][12][5]5.12}{[1][33][]33}}
\newlabel{eje:t5-sidak}{{5.13}{34}{}{ejercicio.5.13}{}}
\newlabel{eje:t5-sidak@cref}{{[ejercicio][13][5]5.13}{[1][34][]34}}
\newlabel{eje:t5-uniforme-p}{{5.14}{34}{}{ejercicio.5.14}{}}
\newlabel{eje:t5-uniforme-p@cref}{{[ejercicio][14][5]5.14}{[1][34][]34}}
\newlabel{eje:t5-dualidad}{{5.15}{34}{}{ejercicio.5.15}{}}
\newlabel{eje:t5-dualidad@cref}{{[ejercicio][15][5]5.15}{[1][34][]34}}
\newlabel{eje:t5-potencia}{{5.16}{34}{}{ejercicio.5.16}{}}
\newlabel{eje:t5-potencia@cref}{{[ejercicio][16][5]5.16}{[1][34][]34}}
\newlabel{ej:t6-recta}{{6.3}{36}{}{ejemplo.6.3}{}}
\newlabel{ej:t6-recta@cref}{{[ejemplo][3][]6.3}{[1][36][]36}}
\newlabel{fig:t6-recta}{{6.1}{36}{Recta de mínimos cuadrados del \cref {ej:t6-recta}. Los segmentos discontinuos son los residuos; la recta pasa por el punto de medias $(3,4)$}{figure.6.1}{}}
\newlabel{fig:t6-recta@cref}{{[figure][1][6]6.1}{[1][36][]36}}
\newlabel{ej:t6-potencia}{{6.9}{38}{}{ejemplo.6.9}{}}
\newlabel{ej:t6-potencia@cref}{{[ejemplo][9][]6.9}{[1][37][]38}}
\newlabel{ej:t6-confusion}{{6.10}{38}{Variable de confusión}{ejemplo.6.10}{}}
\newlabel{ej:t6-confusion@cref}{{[ejemplo][10][]6.10}{[1][38][]38}}
\newlabel{obs:t6-errores}{{6.11}{38}{Errores frecuentes}{observacion.6.11}{}}
\newlabel{obs:t6-errores@cref}{{[observacion][11][]6.11}{[1][38][]38}}
\newlabel{eje:t6-vf}{{6.1}{39}{}{ejercicio.6.1}{}}
\newlabel{eje:t6-vf@cref}{{[ejercicio][1][6]6.1}{[1][39][]39}}
\newlabel{eje:t6-basico}{{6.2}{39}{}{ejercicio.6.2}{}}
\newlabel{eje:t6-basico@cref}{{[ejercicio][2][6]6.2}{[1][39][]39}}
\newlabel{eje:t6-gauss}{{6.3}{39}{}{ejercicio.6.3}{}}
\newlabel{eje:t6-gauss@cref}{{[ejercicio][3][6]6.3}{[1][39][]39}}
\newlabel{eje:t6-examenes}{{6.4}{39}{}{ejercicio.6.4}{}}
\newlabel{eje:t6-examenes@cref}{{[ejercicio][4][6]6.4}{[1][39][]39}}
\newlabel{eje:t6-buceo}{{6.5}{39}{}{ejercicio.6.5}{}}
\newlabel{eje:t6-buceo@cref}{{[ejercicio][5][6]6.5}{[1][39][]39}}
\newlabel{eje:t6-helados}{{6.6}{39}{}{ejercicio.6.6}{}}
\newlabel{eje:t6-helados@cref}{{[ejercicio][6][6]6.6}{[1][39][]39}}
\newlabel{eje:t6-resumen}{{6.7}{39}{}{ejercicio.6.7}{}}
\newlabel{eje:t6-resumen@cref}{{[ejercicio][7][6]6.7}{[1][39][]39}}
\newlabel{eje:t6-loglog}{{6.8}{39}{}{ejercicio.6.8}{}}
\newlabel{eje:t6-loglog@cref}{{[ejercicio][8][6]6.8}{[1][39][]39}}
\newlabel{eje:t6-simpson}{{6.9}{40}{}{ejercicio.6.9}{}}
\newlabel{eje:t6-simpson@cref}{{[ejercicio][9][6]6.9}{[1][39][]40}}
\newlabel{eje:t6-cuadratica}{{6.10}{40}{}{ejercicio.6.10}{}}
\newlabel{eje:t6-cuadratica@cref}{{[ejercicio][10][6]6.10}{[1][40][]40}}
\newlabel{eje:t6-normales}{{6.11}{40}{}{ejercicio.6.11}{}}
\newlabel{eje:t6-normales@cref}{{[ejercicio][11][6]6.11}{[1][40][]40}}
\newlabel{eje:t6-residuos}{{6.12}{40}{}{ejercicio.6.12}{}}
\newlabel{eje:t6-residuos@cref}{{[ejercicio][12][6]6.12}{[1][40][]40}}
\newlabel{eje:t6-origen}{{6.13}{40}{}{ejercicio.6.13}{}}
\newlabel{eje:t6-origen@cref}{{[ejercicio][13][6]6.13}{[1][40][]40}}
\newlabel{eje:t6-transformaciones}{{6.14}{40}{}{ejercicio.6.14}{}}
\newlabel{eje:t6-transformaciones@cref}{{[ejercicio][14][6]6.14}{[1][40][]40}}
\newlabel{eje:t6-regresion-media}{{6.15}{40}{}{ejercicio.6.15}{}}
\newlabel{eje:t6-regresion-media@cref}{{[ejercicio][15][6]6.15}{[1][40][]40}}
\newlabel{ap:soluciones}{{A}{41}{Soluciones de los ejercicios}{appendix.A}{}}
\newlabel{ap:soluciones@cref}{{[appendix][1][2147483647]A}{[1][41][]41}}
\newlabel{sol:eje:t1-vf}{{A.1}{41}{Tema 1}{section*.3}{}}
\newlabel{sol:eje:t1-vf@cref}{{[subappendix][1][2147483647,1]A.1}{[1][41][]41}}
\newlabel{sol:eje:t1-basico}{{A.1}{41}{Tema 1}{section*.4}{}}
\newlabel{sol:eje:t1-basico@cref}{{[subappendix][1][2147483647,1]A.1}{[1][41][]41}}
\newlabel{sol:eje:t1-edades}{{A.1}{42}{Tema 1}{section*.5}{}}
\newlabel{sol:eje:t1-edades@cref}{{[subappendix][1][2147483647,1]A.1}{[1][42][]42}}
\newlabel{sol:eje:t1-cajas}{{A.1}{42}{Tema 1}{section*.6}{}}
\newlabel{sol:eje:t1-cajas@cref}{{[subappendix][1][2147483647,1]A.1}{[1][42][]42}}
\newlabel{sol:eje:t1-agrupados}{{A.1}{43}{Tema 1}{section*.7}{}}
\newlabel{sol:eje:t1-agrupados@cref}{{[subappendix][1][2147483647,1]A.1}{[1][43][]43}}
\newlabel{sol:eje:t1-abreviada}{{A.1}{43}{Tema 1}{section*.8}{}}
\newlabel{sol:eje:t1-abreviada@cref}{{[subappendix][1][2147483647,1]A.1}{[1][43][]43}}
\newlabel{sol:eje:t1-celsius}{{A.1}{43}{Tema 1}{section*.9}{}}
\newlabel{sol:eje:t1-celsius@cref}{{[subappendix][1][2147483647,1]A.1}{[1][43][]43}}
\newlabel{sol:eje:t1-zscore}{{A.1}{44}{Tema 1}{section*.10}{}}
\newlabel{sol:eje:t1-zscore@cref}{{[subappendix][1][2147483647,1]A.1}{[1][44][]44}}
\newlabel{sol:eje:t1-cov-tabla}{{A.1}{44}{Tema 1}{section*.11}{}}
\newlabel{sol:eje:t1-cov-tabla@cref}{{[subappendix][1][2147483647,1]A.1}{[1][44][]44}}
\newlabel{sol:eje:t1-agrupar}{{A.1}{45}{Tema 1}{section*.12}{}}
\newlabel{sol:eje:t1-agrupar@cref}{{[subappendix][1][2147483647,1]A.1}{[1][44][]45}}
\newlabel{sol:eje:t1-mediana-min}{{A.1}{45}{Tema 1}{section*.13}{}}
\newlabel{sol:eje:t1-mediana-min@cref}{{[subappendix][1][2147483647,1]A.1}{[1][45][]45}}
\newlabel{sol:eje:t1-r-invariante}{{A.1}{45}{Tema 1}{section*.14}{}}
\newlabel{sol:eje:t1-r-invariante@cref}{{[subappendix][1][2147483647,1]A.1}{[1][45][]45}}
\newlabel{sol:eje:t1-chebyshev}{{A.1}{46}{Tema 1}{section*.15}{}}
\newlabel{sol:eje:t1-chebyshev@cref}{{[subappendix][1][2147483647,1]A.1}{[1][46][]46}}
\newlabel{sol:eje:t1-welford}{{A.1}{46}{Tema 1}{section*.16}{}}
\newlabel{sol:eje:t1-welford@cref}{{[subappendix][1][2147483647,1]A.1}{[1][46][]46}}
\newlabel{sol:eje:t1-varianza-pares}{{A.1}{47}{Tema 1}{section*.17}{}}
\newlabel{sol:eje:t1-varianza-pares@cref}{{[subappendix][1][2147483647,1]A.1}{[1][46][]47}}
\newlabel{sol:eje:t2-conceptos}{{A.2}{47}{Tema 2}{section*.18}{}}
\newlabel{sol:eje:t2-conceptos@cref}{{[subappendix][2][2147483647,1]A.2}{[1][47][]47}}
\newlabel{sol:eje:t2-enumeracion}{{A.2}{47}{Tema 2}{section*.19}{}}
\newlabel{sol:eje:t2-enumeracion@cref}{{[subappendix][2][2147483647,1]A.2}{[1][47][]47}}
\newlabel{sol:eje:t2-tres-estimadores}{{A.2}{48}{Tema 2}{section*.20}{}}
\newlabel{sol:eje:t2-tres-estimadores@cref}{{[subappendix][2][2147483647,1]A.2}{[1][48][]48}}
\newlabel{sol:eje:t2-poisson}{{A.2}{48}{Tema 2}{section*.21}{}}
\newlabel{sol:eje:t2-poisson@cref}{{[subappendix][2][2147483647,1]A.2}{[1][48][]48}}
\newlabel{sol:eje:t2-expo-sesgo}{{A.2}{49}{Tema 2}{section*.22}{}}
\newlabel{sol:eje:t2-expo-sesgo@cref}{{[subappendix][2][2147483647,1]A.2}{[1][49][]49}}
\newlabel{sol:eje:t2-sumas-solapadas}{{A.2}{49}{Tema 2}{section*.23}{}}
\newlabel{sol:eje:t2-sumas-solapadas@cref}{{[subappendix][2][2147483647,1]A.2}{[1][49][]49}}
\newlabel{sol:eje:t2-benchmark}{{A.2}{49}{Tema 2}{section*.24}{}}
\newlabel{sol:eje:t2-benchmark@cref}{{[subappendix][2][2147483647,1]A.2}{[1][49][]49}}
\newlabel{sol:eje:t2-combinar}{{A.2}{50}{Tema 2}{section*.25}{}}
\newlabel{sol:eje:t2-combinar@cref}{{[subappendix][2][2147483647,1]A.2}{[1][50][]50}}
\newlabel{sol:eje:t2-chebyshev-n}{{A.2}{50}{Tema 2}{section*.26}{}}
\newlabel{sol:eje:t2-chebyshev-n@cref}{{[subappendix][2][2147483647,1]A.2}{[1][50][]50}}
\newlabel{sol:eje:t2-ecm-normal}{{A.2}{50}{Tema 2}{section*.27}{}}
\newlabel{sol:eje:t2-ecm-normal@cref}{{[subappendix][2][2147483647,1]A.2}{[1][50][]50}}
\newlabel{sol:eje:t2-uniforme}{{A.2}{51}{Tema 2}{section*.28}{}}
\newlabel{sol:eje:t2-uniforme@cref}{{[subappendix][2][2147483647,1]A.2}{[1][51][]51}}
\newlabel{sol:eje:t2-max-consistente}{{A.2}{51}{Tema 2}{section*.29}{}}
\newlabel{sol:eje:t2-max-consistente@cref}{{[subappendix][2][2147483647,1]A.2}{[1][51][]51}}
\newlabel{sol:eje:t2-cov-desv}{{A.2}{52}{Tema 2}{section*.30}{}}
\newlabel{sol:eje:t2-cov-desv@cref}{{[subappendix][2][2147483647,1]A.2}{[1][52][]52}}
\newlabel{sol:eje:t2-s-sesgo}{{A.2}{52}{Tema 2}{section*.31}{}}
\newlabel{sol:eje:t2-s-sesgo@cref}{{[subappendix][2][2147483647,1]A.2}{[1][52][]52}}
\newlabel{sol:eje:t2-passk}{{A.2}{52}{Tema 2}{section*.32}{}}
\newlabel{sol:eje:t2-passk@cref}{{[subappendix][2][2147483647,1]A.2}{[1][52][]52}}
\newlabel{sol:eje:t3-vf}{{A.3}{53}{Tema 3}{section*.33}{}}
\newlabel{sol:eje:t3-vf@cref}{{[subappendix][3][2147483647,1]A.3}{[1][53][]53}}
\newlabel{sol:eje:t3-tabla}{{A.3}{53}{Tema 3}{section*.34}{}}
\newlabel{sol:eje:t3-tabla@cref}{{[subappendix][3][2147483647,1]A.3}{[1][53][]53}}
\newlabel{sol:eje:t3-normal-basico}{{A.3}{54}{Tema 3}{section*.35}{}}
\newlabel{sol:eje:t3-normal-basico@cref}{{[subappendix][3][2147483647,1]A.3}{[1][53][]54}}
\newlabel{sol:eje:t3-openstax-media}{{A.3}{54}{Tema 3}{section*.36}{}}
\newlabel{sol:eje:t3-openstax-media@cref}{{[subappendix][3][2147483647,1]A.3}{[1][54][]54}}
\newlabel{sol:eje:t3-openstax-suma}{{A.3}{54}{Tema 3}{section*.37}{}}
\newlabel{sol:eje:t3-openstax-suma@cref}{{[subappendix][3][2147483647,1]A.3}{[1][54][]54}}
\newlabel{sol:eje:t3-elecciones}{{A.3}{54}{Tema 3}{section*.38}{}}
\newlabel{sol:eje:t3-elecciones@cref}{{[subappendix][3][2147483647,1]A.3}{[1][54][]54}}
\newlabel{sol:eje:t3-redondeo}{{A.3}{55}{Tema 3}{section*.39}{}}
\newlabel{sol:eje:t3-redondeo@cref}{{[subappendix][3][2147483647,1]A.3}{[1][55][]55}}
\newlabel{sol:eje:t3-binomial}{{A.3}{55}{Tema 3}{section*.40}{}}
\newlabel{sol:eje:t3-binomial@cref}{{[subappendix][3][2147483647,1]A.3}{[1][55][]55}}
\newlabel{sol:eje:t3-benchmark}{{A.3}{55}{Tema 3}{section*.41}{}}
\newlabel{sol:eje:t3-benchmark@cref}{{[subappendix][3][2147483647,1]A.3}{[1][55][]55}}
\newlabel{sol:eje:t3-tamano}{{A.3}{56}{Tema 3}{section*.42}{}}
\newlabel{sol:eje:t3-tamano@cref}{{[subappendix][3][2147483647,1]A.3}{[1][56][]56}}
\newlabel{sol:eje:t3-t-cuantiles}{{A.3}{56}{Tema 3}{section*.43}{}}
\newlabel{sol:eje:t3-t-cuantiles@cref}{{[subappendix][3][2147483647,1]A.3}{[1][56][]56}}
\newlabel{sol:eje:t3-lineal-normal}{{A.3}{56}{Tema 3}{section*.44}{}}
\newlabel{sol:eje:t3-lineal-normal@cref}{{[subappendix][3][2147483647,1]A.3}{[1][56][]56}}
\newlabel{sol:eje:t3-chi1}{{A.3}{57}{Tema 3}{section*.45}{}}
\newlabel{sol:eje:t3-chi1@cref}{{[subappendix][3][2147483647,1]A.3}{[1][57][]57}}
\newlabel{sol:eje:t3-chi-momentos}{{A.3}{57}{Tema 3}{section*.46}{}}
\newlabel{sol:eje:t3-chi-momentos@cref}{{[subappendix][3][2147483647,1]A.3}{[1][57][]57}}
\newlabel{sol:eje:t3-t-simetrica}{{A.3}{57}{Tema 3}{section*.47}{}}
\newlabel{sol:eje:t3-t-simetrica@cref}{{[subappendix][3][2147483647,1]A.3}{[1][57][]57}}
\newlabel{sol:eje:t4-interpretacion}{{A.4}{58}{Tema 4}{section*.48}{}}
\newlabel{sol:eje:t4-interpretacion@cref}{{[subappendix][4][2147483647,1]A.4}{[1][58][]58}}
\newlabel{sol:eje:t4-momentos-basico}{{A.4}{58}{Tema 4}{section*.49}{}}
\newlabel{sol:eje:t4-momentos-basico@cref}{{[subappendix][4][2147483647,1]A.4}{[1][58][]58}}
\newlabel{sol:eje:t4-geometrica}{{A.4}{59}{Tema 4}{section*.50}{}}
\newlabel{sol:eje:t4-geometrica@cref}{{[subappendix][4][2147483647,1]A.4}{[1][58][]59}}
\newlabel{sol:eje:t4-gamma}{{A.4}{59}{Tema 4}{section*.51}{}}
\newlabel{sol:eje:t4-gamma@cref}{{[subappendix][4][2147483647,1]A.4}{[1][59][]59}}
\newlabel{sol:eje:t4-uniforme-ab}{{A.4}{59}{Tema 4}{section*.52}{}}
\newlabel{sol:eje:t4-uniforme-ab@cref}{{[subappendix][4][2147483647,1]A.4}{[1][59][]59}}
\newlabel{sol:eje:t4-poisson-datos}{{A.4}{60}{Tema 4}{section*.53}{}}
\newlabel{sol:eje:t4-poisson-datos@cref}{{[subappendix][4][2147483647,1]A.4}{[1][60][]60}}
\newlabel{sol:eje:t4-soda}{{A.4}{60}{Tema 4}{section*.54}{}}
\newlabel{sol:eje:t4-soda@cref}{{[subappendix][4][2147483647,1]A.4}{[1][60][]60}}
\newlabel{sol:eje:t4-alturas}{{A.4}{60}{Tema 4}{section*.55}{}}
\newlabel{sol:eje:t4-alturas@cref}{{[subappendix][4][2147483647,1]A.4}{[1][60][]60}}
\newlabel{sol:eje:t4-openstax-z}{{A.4}{61}{Tema 4}{section*.56}{}}
\newlabel{sol:eje:t4-openstax-z@cref}{{[subappendix][4][2147483647,1]A.4}{[1][61][]61}}
\newlabel{sol:eje:t4-proporcion}{{A.4}{61}{Tema 4}{section*.57}{}}
\newlabel{sol:eje:t4-proporcion@cref}{{[subappendix][4][2147483647,1]A.4}{[1][61][]61}}
\newlabel{sol:eje:t4-margen}{{A.4}{61}{Tema 4}{section*.58}{}}
\newlabel{sol:eje:t4-margen@cref}{{[subappendix][4][2147483647,1]A.4}{[1][61][]61}}
\newlabel{sol:eje:t4-emv-poisson}{{A.4}{62}{Tema 4}{section*.59}{}}
\newlabel{sol:eje:t4-emv-poisson@cref}{{[subappendix][4][2147483647,1]A.4}{[1][62][]62}}
\newlabel{sol:eje:t4-emv-expo}{{A.4}{62}{Tema 4}{section*.60}{}}
\newlabel{sol:eje:t4-emv-expo@cref}{{[subappendix][4][2147483647,1]A.4}{[1][62][]62}}
\newlabel{sol:eje:t4-laplace}{{A.4}{62}{Tema 4}{section*.61}{}}
\newlabel{sol:eje:t4-laplace@cref}{{[subappendix][4][2147483647,1]A.4}{[1][62][]62}}
\newlabel{sol:eje:t4-regla3}{{A.4}{63}{Tema 4}{section*.62}{}}
\newlabel{sol:eje:t4-regla3@cref}{{[subappendix][4][2147483647,1]A.4}{[1][63][]63}}
\newlabel{sol:eje:t4-conservador}{{A.4}{63}{Tema 4}{section*.63}{}}
\newlabel{sol:eje:t4-conservador@cref}{{[subappendix][4][2147483647,1]A.4}{[1][63][]63}}
\newlabel{sol:eje:t5-vf}{{A.5}{63}{Tema 5}{section*.64}{}}
\newlabel{sol:eje:t5-vf@cref}{{[subappendix][5][2147483647,1]A.5}{[1][63][]63}}
\newlabel{sol:eje:t5-moneda}{{A.5}{64}{Tema 5}{section*.65}{}}
\newlabel{sol:eje:t5-moneda@cref}{{[subappendix][5][2147483647,1]A.5}{[1][64][]64}}
\newlabel{sol:eje:t5-poligrafo}{{A.5}{64}{Tema 5}{section*.66}{}}
\newlabel{sol:eje:t5-poligrafo@cref}{{[subappendix][5][2147483647,1]A.5}{[1][64][]64}}
\newlabel{sol:eje:t5-radar}{{A.5}{64}{Tema 5}{section*.67}{}}
\newlabel{sol:eje:t5-radar@cref}{{[subappendix][5][2147483647,1]A.5}{[1][64][]64}}
\newlabel{sol:eje:t5-uniforme}{{A.5}{65}{Tema 5}{section*.68}{}}
\newlabel{sol:eje:t5-uniforme@cref}{{[subappendix][5][2147483647,1]A.5}{[1][65][]65}}
\newlabel{sol:eje:t5-t-test}{{A.5}{65}{Tema 5}{section*.69}{}}
\newlabel{sol:eje:t5-t-test@cref}{{[subappendix][5][2147483647,1]A.5}{[1][65][]65}}
\newlabel{sol:eje:t5-openstax-z}{{A.5}{65}{Tema 5}{section*.70}{}}
\newlabel{sol:eje:t5-openstax-z@cref}{{[subappendix][5][2147483647,1]A.5}{[1][65][]65}}
\newlabel{sol:eje:t5-openstax-prop}{{A.5}{66}{Tema 5}{section*.71}{}}
\newlabel{sol:eje:t5-openstax-prop@cref}{{[subappendix][5][2147483647,1]A.5}{[1][66][]66}}
\newlabel{sol:eje:t5-chi-dias}{{A.5}{66}{Tema 5}{section*.72}{}}
\newlabel{sol:eje:t5-chi-dias@cref}{{[subappendix][5][2147483647,1]A.5}{[1][66][]66}}
\newlabel{sol:eje:t5-benford}{{A.5}{66}{Tema 5}{section*.73}{}}
\newlabel{sol:eje:t5-benford@cref}{{[subappendix][5][2147483647,1]A.5}{[1][66][]66}}
\newlabel{sol:eje:t5-dos-modelos}{{A.5}{67}{Tema 5}{section*.74}{}}
\newlabel{sol:eje:t5-dos-modelos@cref}{{[subappendix][5][2147483647,1]A.5}{[1][67][]67}}
\newlabel{sol:eje:t5-multiples}{{A.5}{67}{Tema 5}{section*.75}{}}
\newlabel{sol:eje:t5-multiples@cref}{{[subappendix][5][2147483647,1]A.5}{[1][67][]67}}
\newlabel{sol:eje:t5-sidak}{{A.5}{67}{Tema 5}{section*.76}{}}
\newlabel{sol:eje:t5-sidak@cref}{{[subappendix][5][2147483647,1]A.5}{[1][67][]67}}
\newlabel{sol:eje:t5-uniforme-p}{{A.5}{68}{Tema 5}{section*.77}{}}
\newlabel{sol:eje:t5-uniforme-p@cref}{{[subappendix][5][2147483647,1]A.5}{[1][68][]68}}
\newlabel{sol:eje:t5-dualidad}{{A.5}{68}{Tema 5}{section*.78}{}}
\newlabel{sol:eje:t5-dualidad@cref}{{[subappendix][5][2147483647,1]A.5}{[1][68][]68}}
\newlabel{sol:eje:t5-potencia}{{A.5}{68}{Tema 5}{section*.79}{}}
\newlabel{sol:eje:t5-potencia@cref}{{[subappendix][5][2147483647,1]A.5}{[1][68][]68}}
\newlabel{sol:eje:t6-vf}{{A.6}{69}{Tema 6}{section*.80}{}}
\newlabel{sol:eje:t6-vf@cref}{{[subappendix][6][2147483647,1]A.6}{[1][69][]69}}
\newlabel{sol:eje:t6-basico}{{A.6}{69}{Tema 6}{section*.81}{}}
\newlabel{sol:eje:t6-basico@cref}{{[subappendix][6][2147483647,1]A.6}{[1][69][]69}}
\newlabel{sol:eje:t6-gauss}{{A.6}{70}{Tema 6}{section*.82}{}}
\newlabel{sol:eje:t6-gauss@cref}{{[subappendix][6][2147483647,1]A.6}{[1][70][]70}}
\newlabel{sol:eje:t6-examenes}{{A.6}{70}{Tema 6}{section*.83}{}}
\newlabel{sol:eje:t6-examenes@cref}{{[subappendix][6][2147483647,1]A.6}{[1][70][]70}}
\newlabel{sol:eje:t6-buceo}{{A.6}{70}{Tema 6}{section*.84}{}}
\newlabel{sol:eje:t6-buceo@cref}{{[subappendix][6][2147483647,1]A.6}{[1][70][]70}}
\newlabel{sol:eje:t6-helados}{{A.6}{71}{Tema 6}{section*.85}{}}
\newlabel{sol:eje:t6-helados@cref}{{[subappendix][6][2147483647,1]A.6}{[1][71][]71}}
\newlabel{sol:eje:t6-resumen}{{A.6}{71}{Tema 6}{section*.86}{}}
\newlabel{sol:eje:t6-resumen@cref}{{[subappendix][6][2147483647,1]A.6}{[1][71][]71}}
\newlabel{sol:eje:t6-loglog}{{A.6}{72}{Tema 6}{section*.87}{}}
\newlabel{sol:eje:t6-loglog@cref}{{[subappendix][6][2147483647,1]A.6}{[1][72][]72}}
\newlabel{sol:eje:t6-simpson}{{A.6}{72}{Tema 6}{section*.88}{}}
\newlabel{sol:eje:t6-simpson@cref}{{[subappendix][6][2147483647,1]A.6}{[1][72][]72}}
\newlabel{sol:eje:t6-cuadratica}{{A.6}{72}{Tema 6}{section*.89}{}}
\newlabel{sol:eje:t6-cuadratica@cref}{{[subappendix][6][2147483647,1]A.6}{[1][72][]72}}
\newlabel{sol:eje:t6-normales}{{A.6}{73}{Tema 6}{section*.90}{}}
\newlabel{sol:eje:t6-normales@cref}{{[subappendix][6][2147483647,1]A.6}{[1][73][]73}}
\newlabel{sol:eje:t6-residuos}{{A.6}{73}{Tema 6}{section*.91}{}}
\newlabel{sol:eje:t6-residuos@cref}{{[subappendix][6][2147483647,1]A.6}{[1][73][]73}}
\newlabel{sol:eje:t6-origen}{{A.6}{74}{Tema 6}{section*.92}{}}
\newlabel{sol:eje:t6-origen@cref}{{[subappendix][6][2147483647,1]A.6}{[1][73][]74}}
\newlabel{sol:eje:t6-transformaciones}{{A.6}{74}{Tema 6}{section*.93}{}}
\newlabel{sol:eje:t6-transformaciones@cref}{{[subappendix][6][2147483647,1]A.6}{[1][74][]74}}
\newlabel{sol:eje:t6-regresion-media}{{A.6}{74}{Tema 6}{section*.94}{}}
\newlabel{sol:eje:t6-regresion-media@cref}{{[subappendix][6][2147483647,1]A.6}{[1][74][]74}}
\newlabel{ap:tablas}{{B}{76}{Tablas estadísticas}{appendix.B}{}}
\newlabel{ap:tablas@cref}{{[appendix][2][2147483647]B}{[1][76][]76}}
\newlabel{tab:phi}{{B.1}{76}{Función de distribución $\Phi (z)=\PP (Z\le z)$ de $Z\sim \Nor (0,1)$. La fila da las unidades y la primera cifra decimal de $z$; la columna, la segunda cifra decimal}{table.B.1}{}}
\newlabel{tab:phi@cref}{{[table][1][2147483647,2]B.1}{[1][76][]76}}
\newlabel{tab:zq}{{B.2}{77}{Cuantiles de la normal estándar más usados: $\Phi (z_q)=q$}{table.B.2}{}}
\newlabel{tab:zq@cref}{{[table][2][2147483647,2]B.2}{[1][76][]77}}
\newlabel{tab:t}{{B.3}{78}{Cuantiles $t_{k;q}$ de la distribución $t$ con $k$ grados de libertad: $\PP (T_k\le t_{k;q})=q$. La última fila ($k=\infty $) son los cuantiles de la normal estándar}{table.B.3}{}}
\newlabel{tab:t@cref}{{[table][3][2147483647,2]B.3}{[1][76][]78}}
\newlabel{tab:chi2}{{B.4}{79}{Cuantiles $\chi ^{2}_{k;q}$ de la distribución $\chi ^{2}$ con $k$ grados de libertad: $\PP (\chi ^{2}_k\le \chi ^{2}_{k;q})=q$}{table.B.4}{}}
\newlabel{tab:chi2@cref}{{[table][4][2147483647,2]B.4}{[1][77][]79}}
\makeatother

\begin{document}
\frontmatter
\begin{titlepage}
\centering
\vspace*{2.5cm}
{\large\scshape Anatomía de un LLM}\par
\vspace{0.4cm}
{\normalsize Fundamentos matemáticos de los modelos de lenguaje}\par
\vspace{2.2cm}
\rule{\textwidth}{0.8pt}\par\vspace{0.5cm}
{\small\scshape Fase 0 \textperiodcentered\ Preliminares matemáticos}\par\vspace{0.3cm}
{\Huge\bfseries Módulo 0.9\par}\vspace{0.3cm}
{\LARGE Estadística\par}
\vspace{0.5cm}\rule{\textwidth}{0.4pt}\par
\vspace{1.6cm}
{\large Resumen del módulo}\par
\vfill
{\small Edición de octubre de 2026}\par
\end{titlepage}

\chapter*{Sobre este resumen}
\addcontentsline{toc}{chapter}{Sobre este resumen}
\markboth{Sobre este resumen}{}

Este documento reúne lo indispensable de los apuntes del Módulo 0.9: el texto expositivo, todas las definiciones y todos los resultados (teoremas, proposiciones, lemas y corolarios), las tablas y las observaciones de contenido. Se han suprimido las demostraciones, los ejemplos, las lecturas, las figuras, los ejercicios y sus soluciones. La numeración de temas, secciones, enunciados y ecuaciones es la misma que en los apuntes, de modo que cualquier referencia puede buscarse allí; las que remiten a elementos suprimidos (un ejemplo, una figura, una ecuación de una demostración) conservan su número en los apuntes.

\section*{Notación}
\begin{center}
\small
\begin{tabularx}{\textwidth}{@{}>{\raggedright\arraybackslash}p{5.2cm}X@{}}
\toprule
Símbolo & Significado \\
\midrule
$x_1,\dots,x_n$; $x_{(1)}\le\dots\le x_{(n)}$ & datos observados y datos ordenados \\
$\xbar$, $s^{2}$, $s$ & media, varianza y desviación típica muestrales (observadas) \\
$X_1,\dots,X_n$ \iid & muestra aleatoria simple (variables independientes e idénticamente distribuidas) \\
$\Xbar$, $S^{2}$ & media y varianza muestrales como variables aleatorias \\
$\mu$, $\sigma^{2}$, $p$, $\lambda$, $\theta$ & parámetros de la población \\
$\widehat\theta$ & estimador (o estimación) de $\theta$ \\
$\PP$, $\E$, $\Var$, $\Cov$ & probabilidad, esperanza, varianza y covarianza \\
$\Nor(\mu,\sigma^{2})$, $\Phi$ & distribución normal y función de distribución de la normal estándar \\
$\Ber(p)$, $\Bin(n,p)$, $\Pois(\lambda)$, $\Expo(\lambda)$ & Bernoulli, binomial, Poisson y exponencial de tasa $\lambda$ \\
$z_{q}$, $t_{k;q}$, $\chi^{2}_{k;q}$ & cuantiles de orden $q$ de $\Nor(0,1)$, de $t_k$ y de $\chi^{2}_k$ \\
$H_0$, $H_1$, $\alpha$, $\beta$ & hipótesis nula y alternativa, nivel de significación y probabilidad de error de tipo~II \\
$\ee(\widehat\theta)$, $\ECM(\widehat\theta)$ & error estándar y error cuadrático medio \\
$\convp$ & convergencia en probabilidad \\
\bottomrule
\end{tabularx}
\end{center}
Los cuantiles se escriben siempre por su probabilidad acumulada: $z_{0{,}975}=1{,}96$ es el valor que deja a su izquierda una probabilidad $0{,}975$.


\tableofcontents

\mainmatter
% =====================================================================
\chapter{Estadística descriptiva}
\label{tema:1}
% =====================================================================

Antes de construir modelos conviene mirar los datos: resumirlos con unos pocos números y dibujarlos. Este tema presenta las medidas de centralización (media, mediana), de posición (cuantiles y cuartiles), de dispersión (rango, recorrido intercuartílico, varianza y desviación típica), las representaciones gráficas básicas (histograma y diagrama de caja) y las medidas de relación entre dos variables (covarianza y correlación). Se sigue el capítulo~2 de \emph{Introductory Statistics} de OpenStax~\cite{openstax} (§2.2 a §2.7), el capítulo~10 de Rice~\cite{rice}, «Summarizing Data», y, para la covarianza y la correlación, la lectura~7b del curso 18.05 del MIT~\cite{mit1805} y la §7.3 de Blitzstein y Hwang~\cite{bh}.

\section{Datos, población y muestra}

\begin{definicion}[Población, muestra y variable]\label{def:t1-poblacion}
La \emph{población} es el conjunto de todos los individuos u objetos sobre los que se quiere obtener información; una \emph{muestra} es un subconjunto de la población que efectivamente se observa. Una \emph{variable} es una característica que se mide en cada individuo. La variable es \emph{cuantitativa} si toma valores numéricos con los que tiene sentido operar (discreta si toma valores aislados, como un número de errores; continua si puede tomar cualquier valor de un intervalo, como un tiempo) y \emph{cualitativa} o \emph{categórica} si toma valores en un conjunto de categorías.
\end{definicion}

Un conjunto de datos cuantitativos es una lista $x_1,\dots,x_n$ de números reales, donde $n$ es el \emph{tamaño} de la muestra. Los mismos datos ordenados de menor a mayor se escriben $x_{(1)}\le x_{(2)}\le\dots\le x_{(n)}$; $x_{(1)}$ es el mínimo y $x_{(n)}$, el máximo. En este tema los datos son números fijos y las definiciones son puramente descriptivas; a partir del \cref{tema:2} se interpretarán como valores observados de variables aleatorias.

\addtocounter{definicion}{1}%

\section{Medidas de centralización}

\begin{definicion}[Media, mediana y moda]\label{def:t1-centro}
Sean $x_1,\dots,x_n$ datos reales.
\begin{enumerate}
  \item La \emph{media} (aritmética) es $\displaystyle\xbar=\frac1n\sum_{i=1}^{n}x_i$.
  \item La \emph{mediana} es el valor central de los datos ordenados: $x_{((n+1)/2)}$ si $n$ es impar, y $\tfrac12\bigl(x_{(n/2)}+x_{(n/2+1)}\bigr)$ si $n$ es par.
  \item La \emph{moda} es el valor que aparece con mayor frecuencia (puede no ser única).
\end{enumerate}
\end{definicion}

\addtocounter{definicion}{1}%

El ejemplo ilustra que la media es \emph{sensible a los valores atípicos} y la mediana es \emph{robusta}. Cuando la distribución de los datos es asimétrica con cola larga a la derecha (como las longitudes de respuesta, los tiempos o los ingresos), la media suele ser mayor que la mediana~\cite[§2.6]{openstax}. Las dos medidas tienen una caracterización variacional que explica ese comportamiento.

\begin{proposicion}[Propiedades extremales de la media y la mediana]\label{prop:t1-extremal}
Sean $x_1,\dots,x_n$ datos reales.
\begin{enumerate}
  \item $\sum_{i=1}^{n}(x_i-\xbar)=0$.
  \item Para todo $a\in\R$, $\displaystyle\sum_{i=1}^{n}(x_i-a)^{2}=\sum_{i=1}^{n}(x_i-\xbar)^{2}+n(\xbar-a)^{2}$. En consecuencia, la función $a\mapsto\sum_i(x_i-a)^{2}$ alcanza su mínimo únicamente en $a=\xbar$.
  \item La función $a\mapsto\sum_i\abs{x_i-a}$ alcanza su mínimo en la mediana.
\end{enumerate}
\end{proposicion}

\section{Medidas de posición: cuantiles y diagrama de caja}

La mediana divide los datos ordenados en dos mitades. Los cuartiles los dividen en cuatro partes. Hay varias convenciones para calcularlos que dan resultados ligeramente distintos con pocos datos (los programas estadísticos suelen interpolar); en este módulo se usa la siguiente, que es la de~\cite[§2.3]{openstax} y la habitual en el cálculo a mano.

\begin{definicion}[Cuartiles y recorrido intercuartílico]\label{def:t1-cuartiles}
Sean $x_{(1)}\le\dots\le x_{(n)}$ los datos ordenados. La \emph{mitad inferior} está formada por los $\lfloor n/2\rfloor$ primeros datos y la \emph{mitad superior}, por los $\lfloor n/2\rfloor$ últimos (si $n$ es impar, la mediana no pertenece a ninguna de las dos). El \emph{primer cuartil} $Q_1$ es la mediana de la mitad inferior; el \emph{segundo cuartil} $Q_2$ es la mediana; el \emph{tercer cuartil} $Q_3$ es la mediana de la mitad superior. El \emph{recorrido intercuartílico} es $\RIC=Q_3-Q_1$ y el \emph{rango} es $x_{(n)}-x_{(1)}$.
\end{definicion}

Más en general, el \emph{cuantil} de orden $q\in(0,1)$ (o \emph{percentil} $100q$) es un valor que deja a su izquierda aproximadamente una fracción $q$ de los datos; $Q_1$, $Q_2$ y $Q_3$ son los cuantiles de órdenes $0{,}25$, $0{,}5$ y $0{,}75$.

\begin{definicion}[Resumen de cinco números y diagrama de caja]\label{def:t1-caja}
El \emph{resumen de cinco números} es $(x_{(1)},Q_1,Q_2,Q_3,x_{(n)})$. Un dato se considera \emph{atípico} (criterio de Tukey) si está fuera del intervalo
\[
\bigl[\,Q_1-1{,}5\,\RIC,\;Q_3+1{,}5\,\RIC\,\bigr],
\]
cuyos extremos se llaman \emph{vallas}. El \emph{diagrama de caja} dibuja una caja de $Q_1$ a $Q_3$ con una línea en la mediana y dos «bigotes» que llegan hasta el dato más extremo no atípico de cada lado; los atípicos se marcan por separado~\cite[§2.4]{openstax}.
\end{definicion}

\addtocounter{definicion}{1}%

\addtocounter{figure}{1}%

\begin{definicion}[Histograma]\label{def:t1-histograma}
Se divide un intervalo que contiene los datos en \emph{clases} contiguas $[a_0,a_1),[a_1,a_2),\dots,[a_{k-1},a_k]$ y se cuenta la \emph{frecuencia absoluta} $f_j$ de datos en cada clase. El \emph{histograma} dibuja sobre cada clase un rectángulo de altura proporcional a $f_j$ (si las clases tienen igual amplitud) o a $f_j/(a_j-a_{j-1})$ (si no la tienen, de modo que el \emph{área} sea proporcional a la frecuencia). La \emph{frecuencia relativa} es $f_j/n$.
\end{definicion}

El histograma muestra la forma de la distribución (simetría, asimetría, número de modas), mientras que el diagrama de caja resume la posición y la dispersión y señala los atípicos; la \cref{fig:t1-hist-caja} muestra ambos para un mismo conjunto de datos.

\section{Medidas de dispersión}

\begin{definicion}[Varianza y desviación típica muestrales]\label{def:t1-varianza}
La \emph{varianza muestral} de $x_1,\dots,x_n$ ($n\ge 2$) es
\[
s^{2}=\frac{1}{n-1}\sum_{i=1}^{n}(x_i-\xbar)^{2},
\]
y la \emph{desviación típica muestral} es $s=\sqrt{s^{2}}$. La cantidad $\sum_i(x_i-\xbar)^{2}$ se llama \emph{suma de cuadrados} de los datos.
\end{definicion}

\begin{observacion}[¿Por qué $n-1$?]\label{obs:t1-n-1}
El denominador $n-1$ en lugar de $n$ tiene dos justificaciones. La primera es algebraica: las $n$ desviaciones $x_i-\xbar$ suman cero (\cref{prop:t1-extremal}), de modo que solo $n-1$ de ellas pueden elegirse libremente; se dice que la suma de cuadrados tiene $n-1$ \emph{grados de libertad}. La segunda es estadística: si los datos son una muestra aleatoria de una población con varianza $\sigma^{2}$, la versión con $n-1$ acierta «en promedio», es decir, es un estimador insesgado de $\sigma^{2}$, mientras que la versión con $n$ la subestima sistemáticamente. Esto se demuestra en el \cref{tema:2} (\cref{teo:t2-s2}). Con denominador $n$ se obtiene la \emph{varianza poblacional} de los propios datos, $\frac1n\sum_i(x_i-\xbar)^{2}$, que es la que se usa cuando los datos \emph{son} toda la población o cuando se normaliza un vector, como en la normalización por capas de una red (módulo~11).
\end{observacion}

\begin{proposicion}[Fórmula de cálculo]\label{prop:t1-abreviada}
$\displaystyle\sum_{i=1}^{n}(x_i-\xbar)^{2}=\sum_{i=1}^{n}x_i^{2}-n\xbar^{2}$, y por tanto $s^{2}=\dfrac{\sum_i x_i^{2}-n\xbar^{2}}{n-1}$.
\end{proposicion}

\addtocounter{definicion}{1}%

\begin{proposicion}[Cambios de escala y de origen]\label{prop:t1-lineal}
Si $y_i=ax_i+b$ con $a,b\in\R$, entonces $\ybar=a\xbar+b$, $s_y^{2}=a^{2}s_x^{2}$ y $s_y=\abs{a}\,s_x$. Si $a>0$, la mediana y los cuartiles se transforman igual que la media y $\RIC_y=a\,\RIC_x$.
\end{proposicion}

\begin{definicion}[Tipificación]\label{def:t1-z}
Si $s>0$, la \emph{puntuación tipificada} (o puntuación $z$) del dato $x_i$ es $z_i=(x_i-\xbar)/s$. Por la \cref{prop:t1-lineal}, los datos tipificados tienen media $0$ y desviación típica $1$; $z_i$ mide a cuántas desviaciones típicas está $x_i$ de la media.
\end{definicion}

\section{Datos bivariantes: covarianza y correlación}

Cuando en cada individuo se miden dos variables se obtienen pares $(x_1,y_1),\dots,(x_n,y_n)$, que se representan en un \emph{diagrama de dispersión}.

\begin{definicion}[Covarianza y correlación muestrales]\label{def:t1-cov}
La \emph{covarianza muestral} y el \emph{coeficiente de correlación} (de Pearson) son
\[
s_{xy}=\frac{1}{n-1}\sum_{i=1}^{n}(x_i-\xbar)(y_i-\ybar),
\qquad
r=\frac{s_{xy}}{s_x\,s_y}=\frac{\sum_i(x_i-\xbar)(y_i-\ybar)}{\sqrt{\sum_i(x_i-\xbar)^{2}}\,\sqrt{\sum_i(y_i-\ybar)^{2}}},
\]
donde se supone $s_x,s_y>0$.
\end{definicion}

El signo de cada producto $(x_i-\xbar)(y_i-\ybar)$ es positivo cuando el punto está en el cuadrante superior derecho o inferior izquierdo respecto de $(\xbar,\ybar)$ y negativo en los otros dos; $s_{xy}>0$ indica que predominan los primeros, es decir, una tendencia creciente. La covarianza depende de las unidades; la correlación no (\cref{eje:t1-r-invariante}).

\begin{teorema}[Acotación de la correlación]\label{teo:t1-r}
Se cumple $-1\le r\le 1$. Además, $\abs{r}=1$ si y solo si existen $a\neq 0$ y $b$ tales que $y_i=ax_i+b$ para todo $i$, y entonces el signo de $r$ es el de $a$.
\end{teorema}

\addtocounter{definicion}{1}%

Los mismos conceptos existen para variables aleatorias. Si $(X,Y)$ es un vector aleatorio con varianzas finitas y positivas, su \emph{covarianza} y su \emph{correlación} son
\[
\Cov(X,Y)=\E\bigl[(X-\E X)(Y-\E Y)\bigr]=\E[XY]-\E[X]\E[Y],
\qquad
\Cor(X,Y)=\frac{\Cov(X,Y)}{\sqrt{\Var X\,\Var Y}},
\]
y también $\abs{\Cor(X,Y)}\le1$~\cite[§7.3]{bh}. Las fórmulas muestrales de la \cref{def:t1-cov} son las análogas con los datos (salvo el factor $n/(n-1)$, que se cancela en $r$).

\addtocounter{definicion}{1}%

\addtocounter{definicion}{1}%


% =====================================================================
\chapter{Muestreo y estimadores}
\label{tema:2}
% =====================================================================

La estadística inferencial trata los datos como valores observados de variables aleatorias cuya distribución depende de parámetros desconocidos. Este tema introduce el modelo de muestra aleatoria simple, la noción de estimador y los criterios para juzgar su calidad: sesgo, varianza, error estándar, error cuadrático medio y consistencia. Se demuestran los dos resultados centrales del muestreo: la media muestral es un estimador insesgado de la media y la varianza muestral con denominador $n-1$ lo es de la varianza. Se sigue el capítulo~6 de Wasserman~\cite{wasserman}, «Models, Statistical Inference and Learning», la §6.3 de Blitzstein y Hwang~\cite{bh}, «Sample moments», la clase~12 del curso Stats~200 de Stanford~\cite{stats200} y la lectura~10a del curso 18.05 del MIT~\cite{mit1805}.

\section{Recordatorio: esperanza y varianza}

Se recogen las propiedades de la esperanza y la varianza de variables aleatorias (discretas o continuas) que se usarán continuamente~\cite[cap.~4 y §7.3]{bh}. Todas las variables que aparecen se suponen con varianza finita.

\begin{proposicion}[Propiedades de la esperanza y la varianza]\label{prop:t2-props}
Sean $X,Y,X_1,\dots,X_n$ variables aleatorias y $a,b\in\R$.
\begin{enumerate}
  \item Linealidad: $\E[aX+bY]=a\E[X]+b\E[Y]$ y, en general, $\E\bigl[\sum_iX_i\bigr]=\sum_i\E[X_i]$.
  \item $\Var(X)=\E[(X-\E X)^{2}]=\E[X^{2}]-(\E X)^{2}$, y $\Var(aX+b)=a^{2}\Var(X)$.
  \item $\Var\bigl(\sum_iX_i\bigr)=\sum_i\Var(X_i)+2\sum_{i<j}\Cov(X_i,X_j)$.
  \item Si $X$ e $Y$ son independientes, $\E[XY]=\E[X]\E[Y]$ y, por tanto, $\Cov(X,Y)=0$. En consecuencia, si $X_1,\dots,X_n$ son independientes, $\Var\bigl(\sum_iX_i\bigr)=\sum_i\Var(X_i)$.
  \item La covarianza es bilineal,
  \[
  \Cov\Bigl(\sum_ia_iX_i,\sum_jb_jY_j\Bigr)=\sum_{i,j}a_ib_j\Cov(X_i,Y_j),
  \]
  y $\Cov(X,X)=\Var(X)$.
\end{enumerate}
\end{proposicion}

\section{Muestra aleatoria, parámetro y estadístico}

\begin{definicion}[Muestra aleatoria simple]\label{def:t2-mas}
Una \emph{muestra aleatoria simple} de tamaño $n$ de una distribución $F$ es una colección de variables aleatorias $X_1,\dots,X_n$ independientes y todas con distribución $F$; se dice que son \emph{independientes e idénticamente distribuidas} (\iid) y se escribe $X_1,\dots,X_n\sim F$ \iid\ Los datos $x_1,\dots,x_n$ son una realización (valores observados) de la muestra.
\end{definicion}

El modelo describe el muestreo con reemplazamiento de una población, o el muestreo sin reemplazamiento cuando la población es mucho mayor que la muestra, y también la repetición independiente de un experimento: lanzar una moneda, generar una respuesta de un modelo de lenguaje con una semilla distinta o elegir al azar una pregunta de un banco muy grande.

\begin{definicion}[Modelo paramétrico, parámetro y estadístico]\label{def:t2-parametro}
Un \emph{modelo paramétrico} es una familia de distribuciones $\{F_\theta:\theta\in\Theta\}$ indexada por un \emph{parámetro} $\theta$ (un número o un vector) que toma valores en el \emph{espacio paramétrico} $\Theta$. Un \emph{estadístico} es cualquier función $T=g(X_1,\dots,X_n)$ de la muestra que no depende de parámetros desconocidos; es, por tanto, una variable aleatoria, y su distribución se llama \emph{distribución muestral} de $T$.
\end{definicion}

Por ejemplo, $\{\Ber(p):p\in[0,1]\}$, $\{\Pois(\lambda):\lambda>0\}$ y $\{\Nor(\mu,\sigma^{2}):\mu\in\R,\sigma>0\}$ son modelos paramétricos. Los estadísticos más importantes son la \emph{media muestral} y la \emph{varianza muestral}
\[
\Xbar=\frac1n\sum_{i=1}^{n}X_i,\qquad S^{2}=\frac{1}{n-1}\sum_{i=1}^{n}(X_i-\Xbar)^{2},
\]
que son las cantidades del \cref{tema:1} vistas como variables aleatorias. Se reserva la mayúscula para la variable aleatoria y la minúscula ($\xbar$, $s^{2}$) para su valor observado.

\section{Estimadores y sus propiedades}

\begin{definicion}[Estimador, sesgo, error estándar y error cuadrático medio]\label{def:t2-estimador}
Un \emph{estimador} de un parámetro $\theta$ es un estadístico $\widehat\theta=\widehat\theta(X_1,\dots,X_n)$ que se usa para aproximar $\theta$; su valor con unos datos concretos es una \emph{estimación}. Se definen:
\begin{enumerate}
  \item el \emph{sesgo}: $\sesgo(\widehat\theta)=\E_\theta[\widehat\theta]-\theta$; el estimador es \emph{insesgado} si su sesgo es nulo para todo $\theta\in\Theta$;
  \item el \emph{error estándar}: $\ee(\widehat\theta)=\sqrt{\Var_\theta(\widehat\theta)}$;
  \item el \emph{error cuadrático medio}: $\ECM(\widehat\theta)=\E_\theta\bigl[(\widehat\theta-\theta)^{2}\bigr]$.
\end{enumerate}
El subíndice $\theta$ indica que las esperanzas se calculan con la distribución $F_\theta$.
\end{definicion}

El sesgo mide el error sistemático y el error estándar, la dispersión del estimador alrededor de su propia media. El error estándar suele depender de $\theta$, que es desconocido; al sustituir $\theta$ por su estimación se obtiene el \emph{error estándar estimado}, que se escribe $\widehat{\ee}$.

\begin{proposicion}[Descomposición sesgo--varianza]\label{prop:t2-ecm}
$\ECM(\widehat\theta)=\sesgo(\widehat\theta)^{2}+\Var(\widehat\theta)$.
\end{proposicion}

Un estimador ligeramente sesgado puede tener menor $\ECM$ que uno insesgado si su varianza es bastante menor (\cref{eje:t2-ecm-normal}); este compromiso entre sesgo y varianza reaparece en el aprendizaje automático (módulo~8).

\section{La media muestral}

\begin{teorema}[Media y varianza de la media muestral]\label{teo:t2-media}
Sean $X_1,\dots,X_n$ \iid\ con media $\mu$ y varianza $\sigma^{2}$. Entonces
\[
\E[\Xbar]=\mu,\qquad \Var(\Xbar)=\frac{\sigma^{2}}{n},\qquad \ee(\Xbar)=\frac{\sigma}{\sqrt n}.
\]
En particular, $\Xbar$ es un estimador insesgado de $\mu$.
\end{teorema}

El error estándar de la media decrece como $1/\sqrt n$: para dividirlo por dos hay que multiplicar el tamaño de la muestra por cuatro. Esta ley de la raíz cuadrada gobierna el coste de casi toda estimación por muestreo, desde las encuestas hasta la evaluación de modelos.

\addtocounter{definicion}{1}%

\section{La varianza muestral}

\begin{teorema}[Insesgadez de la varianza muestral]\label{teo:t2-s2}
Sean $X_1,\dots,X_n$ \iid\ con media $\mu$ y varianza $\sigma^{2}$, con $n\ge2$. Entonces $\E[S^{2}]=\sigma^{2}$.
\end{teorema}

La demostración explica el $n-1$: las desviaciones se miden respecto de $\Xbar$, que está más cerca de los datos que la verdadera media $\mu$ (es la que minimiza la suma de cuadrados), de modo que $\sum_i(X_i-\Xbar)^{2}$ es en promedio menor que $\sum_i(X_i-\mu)^{2}$, exactamente en una cantidad $\sigma^{2}$.

\begin{corolario}[Sesgo de la varianza con denominador $n$]\label{cor:t2-sesgo}
El estimador $\widehat\sigma^{2}=\frac1n\sum_i(X_i-\Xbar)^{2}=\frac{n-1}{n}S^{2}$ cumple $\E[\widehat\sigma^{2}]=\frac{n-1}{n}\sigma^{2}$, luego tiene sesgo $-\sigma^{2}/n$: subestima la varianza, aunque el sesgo tiende a cero cuando $n\to\infty$.
\end{corolario}

\begin{observacion}
Que $S^{2}$ sea insesgado para $\sigma^{2}$ no implica que $S$ lo sea para $\sigma$: en general $\E[S]<\sigma$ (\cref{eje:t2-s-sesgo}). La insesgadez no se conserva por transformaciones no lineales.
\end{observacion}

\section{Consistencia y ley de los grandes números}

Un estimador razonable debe acercarse al parámetro cuando crece el tamaño de la muestra. Se escribe $\widehat\theta_n$ para subrayar la dependencia de $n$.

\begin{definicion}[Consistencia]\label{def:t2-consistencia}
Una sucesión de variables aleatorias $Y_n$ \emph{converge en probabilidad} a una constante $c$, y se escribe $Y_n\convp c$, si para todo $\varepsilon>0$
\[
\lim_{n\to\infty}\PP\bigl(\abs{Y_n-c}>\varepsilon\bigr)=0.
\]
Un estimador $\widehat\theta_n$ es \emph{consistente} si $\widehat\theta_n\convp\theta$ para todo $\theta\in\Theta$.
\end{definicion}

\begin{lema}[Desigualdades de Markov y de Chebyshev]\label{lem:t2-chebyshev}
Si $Y\ge0$ y $a>0$, entonces $\PP(Y\ge a)\le\E[Y]/a$. En consecuencia, si $W$ tiene media $\mu_W$ y varianza finita, para todo $\varepsilon>0$
\[
\PP(\abs{W-\mu_W}\ge\varepsilon)\le\frac{\Var(W)}{\varepsilon^{2}}.
\]
\end{lema}

\begin{proposicion}[Un $\ECM$ que tiende a cero implica consistencia]\label{prop:t2-ecm-consistencia}
Si $\ECM(\widehat\theta_n)\to0$ para todo $\theta$, entonces $\widehat\theta_n$ es consistente. En particular, basta que el sesgo y la varianza tiendan a cero.
\end{proposicion}

\begin{teorema}[Ley débil de los grandes números]\label{teo:t2-lgn}
Si $X_1,X_2,\dots$ son \iid\ con media $\mu$ y varianza finita $\sigma^{2}$, entonces $\Xbar_n\convp\mu$; es decir, la media muestral es un estimador consistente de la media.
\end{teorema}

La cota de Chebyshev es válida para cualquier distribución, pero es muy conservadora; en el \cref{tema:3} el teorema central del límite dará aproximaciones mucho más precisas.

\addtocounter{definicion}{1}%

\addtocounter{definicion}{1}%


% =====================================================================
\chapter{Distribuciones muestrales y teorema central del límite}
\label{tema:3}
% =====================================================================

Un estadístico es una variable aleatoria y, para hacer inferencia, hace falta conocer (al menos aproximadamente) su distribución. Este tema estudia la distribución de la media muestral: exacta cuando los datos son normales y aproximada, gracias al teorema central del límite, en cualquier otro caso. Se incluyen la aproximación normal a la binomial y las distribuciones $\chi^{2}$ y $t$ de Student, que se usarán en los intervalos de confianza y en los contrastes. Se sigue el capítulo~10 de Blitzstein y Hwang~\cite{bh} (§10.3 «Central limit theorem» y §10.4 «Chi-Square and Student-t»), la lectura~6b del curso 18.05 del MIT~\cite{mit1805}, «Central Limit Theorem and the Law of Large Numbers», y el capítulo~7 de OpenStax~\cite{openstax} (§7.1 y §7.2).

\section{La distribución normal}

\begin{definicion}[Distribución normal]\label{def:t3-normal}
Una variable $X$ tiene distribución \emph{normal} de media $\mu\in\R$ y varianza $\sigma^{2}>0$, y se escribe $X\sim\Nor(\mu,\sigma^{2})$, si su densidad es
\[
f(x)=\frac{1}{\sigma\sqrt{2\pi}}\exp\Bigl(-\frac{(x-\mu)^{2}}{2\sigma^{2}}\Bigr),\qquad x\in\R.
\]
La \emph{normal estándar} es $Z\sim\Nor(0,1)$; su función de distribución se denota $\Phi(z)=\PP(Z\le z)$ y su cuantil de orden $q$, $z_q=\Phi^{-1}(q)$.
\end{definicion}

La función $\Phi$ no tiene expresión mediante funciones elementales y se tabula (\cref{tab:phi}). Por la simetría de la densidad respecto de~$0$, $\Phi(-z)=1-\Phi(z)$ y $z_{1-q}=-z_q$.

\begin{proposicion}[Tipificación]\label{prop:t3-tipificacion}
Si $X\sim\Nor(\mu,\sigma^{2})$ y $a\neq0$, entonces $aX+b\sim\Nor(a\mu+b,a^{2}\sigma^{2})$. En particular, $Z=(X-\mu)/\sigma\sim\Nor(0,1)$ y
\[
\PP(X\le x)=\Phi\Bigl(\frac{x-\mu}{\sigma}\Bigr).
\]
\end{proposicion}

La demostración es el \cref{eje:t3-lineal-normal}. Con la tabla se obtiene la conocida regla empírica: una normal toma valores a menos de $1$, $2$ y $3$ desviaciones típicas de su media con probabilidades $2\Phi(1)-1\approx0{,}683$, $2\Phi(2)-1\approx0{,}954$ y $2\Phi(3)-1\approx0{,}997$.

\begin{teorema}[Combinaciones lineales de normales independientes]\label{teo:t3-suma-normales}
Si $X_1,\dots,X_n$ son independientes con $X_i\sim\Nor(\mu_i,\sigma_i^{2})$, entonces $\sum_ia_iX_i\sim\Nor\bigl(\sum_ia_i\mu_i,\sum_ia_i^{2}\sigma_i^{2}\bigr)$ para cualesquiera $a_i$ no todos nulos. En particular, si $X_1,\dots,X_n\sim\Nor(\mu,\sigma^{2})$ son \iid,
\[
\Xbar\sim\Nor\Bigl(\mu,\frac{\sigma^{2}}{n}\Bigr).
\]
\end{teorema}

La media y la varianza se siguen de la \cref{prop:t2-props}; que la distribución siga siendo normal se demuestra con funciones generadoras de momentos~\cite[§6.6]{bh} y se admite aquí sin demostración.

\section{El teorema central del límite}

Cuando los datos no son normales, la distribución exacta de $\Xbar$ suele ser complicada. El teorema central del límite afirma que, sin embargo, para $n$ grande es aproximadamente normal, sea cual sea la distribución de partida, con tal de que tenga varianza finita.

\begin{teorema}[Teorema central del límite]\label{teo:t3-tcl}
Sean $X_1,X_2,\dots$ \iid\ con media $\mu$ y varianza $0<\sigma^{2}<\infty$, y sea $S_n=X_1+\dots+X_n$. Entonces, para todo $z\in\R$,
\[
\lim_{n\to\infty}\PP\Bigl(\frac{\Xbar_n-\mu}{\sigma/\sqrt n}\le z\Bigr)=\lim_{n\to\infty}\PP\Bigl(\frac{S_n-n\mu}{\sigma\sqrt n}\le z\Bigr)=\Phi(z).
\]
\end{teorema}

Las dos variables tipificadas son la misma (basta multiplicar numerador y denominador por $n$). La demostración, con funciones generadoras de momentos, está en~\cite[§10.3]{bh}; aquí interesa su uso.

\begin{observacion}[Uso práctico]\label{obs:t3-uso}
Para $n$ «grande» se aproxima
\[
\Xbar\approx\Nor\Bigl(\mu,\frac{\sigma^{2}}{n}\Bigr),\qquad S_n\approx\Nor(n\mu,n\sigma^{2}).
\]
Una regla habitual es $n\ge30$~\cite[§7.1]{openstax}, pero la calidad de la aproximación depende de la forma de la distribución: con datos simétricos basta un $n$ pequeño y con datos muy asimétricos hace falta bastante más. El teorema exige varianza finita: para distribuciones de cola muy pesada, como la de Cauchy (o la $t$ con un grado de libertad), la media muestral no se estabiliza~\cite[§10.3]{bh}.
\end{observacion}

\addtocounter{definicion}{1}%

\section{Aproximación normal a la binomial}

Si $X\sim\Bin(n,p)$, entonces $X$ tiene la misma distribución que la suma de $n$ variables $\Ber(p)$ \iid, con $\mu=p$ y $\sigma^{2}=p(1-p)$. El teorema central del límite da el resultado clásico de De Moivre y Laplace.

\begin{corolario}[Aproximación normal a la binomial]\label{cor:t3-binomial}
Si $X\sim\Bin(n,p)$ con $0<p<1$, para $n$ grande
\[
X\approx\Nor\bigl(np,\,np(1-p)\bigr),\qquad \widehat p=\frac Xn\approx\Nor\Bigl(p,\,\frac{p(1-p)}{n}\Bigr).
\]
Una regla habitual para usarla es $np\ge10$ y $n(1-p)\ge10$.
\end{corolario}

\begin{definicion}[Corrección por continuidad]\label{def:t3-cc}
Al aproximar una variable con valores enteros por una continua, cada valor entero $k$ se identifica con el intervalo $[k-\frac12,k+\frac12]$. Así, para enteros $a\le b$,
\[
\PP(a\le X\le b)\approx\Phi\Bigl(\frac{b+\frac12-np}{\sqrt{np(1-p)}}\Bigr)-\Phi\Bigl(\frac{a-\frac12-np}{\sqrt{np(1-p)}}\Bigr).
\]
\end{definicion}

\addtocounter{definicion}{1}%

\section{Distribuciones \texorpdfstring{$\chi^{2}$}{chi cuadrado} y \texorpdfstring{$t$}{t} de Student}

Cuando $\sigma$ es desconocida y se sustituye por $S$, la variable $(\Xbar-\mu)/(S/\sqrt n)$ ya no es normal. Su distribución exacta, para datos normales, se describe con dos familias derivadas de la normal~\cite[§10.4]{bh}.

\begin{definicion}[Distribuciones $\chi^{2}$ y $t$]\label{def:t3-chi-t}
Sean $Z_1,\dots,Z_k\sim\Nor(0,1)$ \iid\ La distribución de $V=Z_1^{2}+\dots+Z_k^{2}$ se llama \emph{$\chi^{2}$ con $k$ grados de libertad} y se escribe $V\sim\chi^{2}_k$. Si $Z\sim\Nor(0,1)$ y $V\sim\chi^{2}_k$ son independientes, la distribución de
\[
T=\frac{Z}{\sqrt{V/k}}
\]
se llama \emph{$t$ de Student con $k$ grados de libertad} y se escribe $T\sim t_k$. Sus cuantiles se denotan $\chi^{2}_{k;q}$ y $t_{k;q}$ (\cref{tab:t,tab:chi2}).
\end{definicion}

La $\chi^{2}_k$ toma valores positivos, tiene media $k$ y varianza $2k$ (\cref{eje:t3-chi-momentos}). La $t_k$ es simétrica respecto de~$0$ (\cref{eje:t3-t-simetrica}), tiene colas más pesadas que la normal y converge a ella cuando $k\to\infty$: por ejemplo, $t_{5;0{,}975}=2{,}571$, $t_{30;0{,}975}=2{,}042$ y $z_{0{,}975}=1{,}960$.

\begin{teorema}[Teorema de Fisher]\label{teo:t3-fisher}
Sean $X_1,\dots,X_n\sim\Nor(\mu,\sigma^{2})$ \iid\ con $n\ge2$. Entonces $\Xbar$ y $S^{2}$ son independientes, $\Xbar\sim\Nor(\mu,\sigma^{2}/n)$ y
\[
\frac{(n-1)S^{2}}{\sigma^{2}}\sim\chi^{2}_{n-1}.
\]
\end{teorema}

Se admite sin demostración~\cite[§10.4]{bh}. Los $n-1$ grados de libertad son los de la \cref{obs:t1-n-1}: la suma de cuadrados respecto de $\Xbar$ equivale a una suma de $n-1$ cuadrados de normales estándar independientes. Como $\E[\chi^{2}_{n-1}]=n-1$, el teorema reproduce, para datos normales, el resultado $\E[S^{2}]=\sigma^{2}$ del \cref{teo:t2-s2}, y como $\Var(\chi^{2}_{n-1})=2(n-1)$, da $\Var(S^{2})=2\sigma^{4}/(n-1)$, que se usó en el \cref{eje:t2-ecm-normal}.

\begin{corolario}[Estadístico $t$]\label{cor:t3-t}
En las condiciones del \cref{teo:t3-fisher},
\[
T=\frac{\Xbar-\mu}{S/\sqrt n}\sim t_{n-1}.
\]
\end{corolario}

\addtocounter{definicion}{1}%

\addtocounter{definicion}{1}%


% =====================================================================
\chapter{Estimación puntual e intervalos de confianza}
\label{tema:4}
% =====================================================================

El \cref{tema:2} enseñó a juzgar un estimador; este tema enseña a construirlo. Se presentan los dos métodos generales de estimación puntual, el método de los momentos y el de máxima verosimilitud, y se aplican a los modelos de Bernoulli, Poisson, normal y exponencial. Después se pasa de la estimación puntual a la estimación por intervalos: intervalos de confianza para la media con varianza conocida y desconocida, para proporciones y cálculo del tamaño muestral. Se sigue el capítulo~9 de Wasserman~\cite{wasserman}, «Parametric Inference» (método de los momentos y máxima verosimilitud), las clases~12, 13 y~18 de Stats~200~\cite{stats200}, las lecturas~10b, 22 y~23a del curso 18.05 del MIT~\cite{mit1805} y el capítulo~8 de OpenStax~\cite{openstax}.

\section{El método de los momentos}

Sea $X_1,\dots,X_n$ una muestra \iid\ de un modelo $F_\theta$ con $\theta=(\theta_1,\dots,\theta_k)$. El \emph{momento poblacional} de orden $j$ es $m_j(\theta)=\E_\theta[X^{j}]$ y el \emph{momento muestral} de orden $j$ es $\widehat m_j=\frac1n\sum_iX_i^{j}$. Por la ley de los grandes números, $\widehat m_j\convp m_j(\theta)$.

\begin{definicion}[Estimador por el método de los momentos]\label{def:t4-momentos}
El \emph{estimador por el método de los momentos} $\widehat\theta$ es la solución del sistema de $k$ ecuaciones
\[
m_j(\widehat\theta)=\widehat m_j,\qquad j=1,\dots,k,
\]
es decir, el valor del parámetro con el que los primeros $k$ momentos del modelo coinciden con los de la muestra~\cite[cap.~9]{wasserman}.
\end{definicion}

\addtocounter{definicion}{1}%

El método es sencillo y produce estimadores consistentes en condiciones generales, pero no siempre los mejores y, a veces, incompatibles con los datos (\cref{eje:t4-uniforme-ab}).

\section{Máxima verosimilitud}

\begin{definicion}[Verosimilitud y estimador de máxima verosimilitud]\label{def:t4-emv}
Sean $x_1,\dots,x_n$ los datos observados y $f(x\mid\theta)$ la función de probabilidad (caso discreto) o la densidad (caso continuo) del modelo. La \emph{función de verosimilitud} y la \emph{log-verosimilitud} son
\[
\Lik(\theta)=\prod_{i=1}^{n}f(x_i\mid\theta),\qquad \ell(\theta)=\ln\Lik(\theta)=\sum_{i=1}^{n}\ln f(x_i\mid\theta).
\]
El \emph{estimador de máxima verosimilitud} (EMV) es el valor $\widehat\theta=\argmax_{\theta\in\Theta}\Lik(\theta)$~\cite[cap.~9]{wasserman}.
\end{definicion}

La verosimilitud es la probabilidad (o densidad) de los datos observados vista como función del parámetro, con los datos fijos. Como el logaritmo es estrictamente creciente, maximizar $\Lik$ equivale a maximizar $\ell$, que es una suma y se deriva con más facilidad. El procedimiento habitual es resolver la \emph{ecuación de verosimilitud} $\ell'(\theta)=0$ y comprobar que la solución es un máximo global.

\begin{teorema}[EMV de Bernoulli]\label{teo:t4-emv-bernoulli}
Si $x_1,\dots,x_n\in\{0,1\}$ son observaciones \iid\ de $\Ber(p)$ y $k=\sum_ix_i$ es el número de unos, con $0<k<n$, el EMV es $\widehat p=k/n=\xbar$.
\end{teorema}

Si $k=0$ o $k=n$, el máximo está en $\widehat p=0$ o $\widehat p=1$, respectivamente, y la fórmula $\widehat p=k/n$ sigue siendo válida.

\begin{teorema}[EMV de la normal]\label{teo:t4-emv-normal}
Si $x_1,\dots,x_n$ son observaciones \iid\ de $\Nor(\mu,\sigma^{2})$, no todas iguales, el EMV de $(\mu,\sigma^{2})$ es
\[
\widehat\mu=\xbar,\qquad\widehat\sigma^{2}=\frac1n\sum_{i=1}^{n}(x_i-\xbar)^{2}.
\]
\end{teorema}

El EMV de $\sigma^{2}$ coincide con el estimador por momentos y es sesgado (\cref{cor:t2-sesgo}); por eso, para datos, se usa casi siempre $S^{2}$.

\addtocounter{definicion}{1}%

\begin{observacion}[Propiedades del EMV]\label{obs:t4-propiedades}
Se enuncian sin demostración, como ampliación~\cite[cap.~9]{wasserman}: (a) \emph{invariancia}: si $\widehat\theta$ es el EMV de $\theta$, el EMV de $g(\theta)$ es $g(\widehat\theta)$; por ejemplo, el EMV de $\sigma$ es $\sqrt{\widehat\sigma^{2}}$; (b) en condiciones de regularidad (que no se cumplen en el \cref{ej:t4-uniforme}), el EMV es consistente y aproximadamente normal para $n$ grande, con un error estándar que se puede estimar a partir de la curvatura de $\ell$ en el máximo.
\end{observacion}

\section{Intervalos de confianza}

Una estimación puntual no informa de su precisión. Un intervalo de confianza acompaña la estimación de un margen de error con una garantía probabilística.

\begin{definicion}[Intervalo de confianza]\label{def:t4-ic}
Sea $0<\alpha<1$. Un \emph{intervalo de confianza} de nivel $1-\alpha$ para $\theta$ es un intervalo $[L,U]$ cuyos extremos $L=L(X_1,\dots,X_n)$ y $U=U(X_1,\dots,X_n)$ son estadísticos y tal que
\[
\PP_\theta(L\le\theta\le U)\ge1-\alpha\qquad\text{para todo }\theta\in\Theta.
\]
La probabilidad $\PP_\theta(L\le\theta\le U)$ se llama \emph{cobertura}~\cite[cap.~6]{wasserman}.
\end{definicion}

\begin{observacion}[Interpretación]\label{obs:t4-interpretacion}
Lo aleatorio es el intervalo, no el parámetro. Si se repitiera muchas veces el experimento y se calculara cada vez el intervalo de nivel $95\,\%$, aproximadamente el $95\,\%$ de esos intervalos contendría el verdadero $\theta$. Una vez observados los datos, el intervalo concreto (por ejemplo, $[0{,}79;\allowbreak\,0{,}84]$) contiene o no contiene a $\theta$; no es correcto decir que $\theta$ está en él «con probabilidad $0{,}95$» en el sentido frecuentista~\cite[lectura~23a]{mit1805}.
\end{observacion}

El método general para construir intervalos usa un \emph{pivote}: una función $Q(X_1,\dots,X_n;\theta)$ cuya distribución no depende de $\theta$.

\begin{teorema}[Intervalo para la media con $\sigma$ conocida]\label{teo:t4-ic-z}
Sean $X_1,\dots,X_n\sim\Nor(\mu,\sigma^{2})$ \iid\ con $\sigma$ conocida. El intervalo
\[
\Bigl[\Xbar-z_{1-\alpha/2}\frac{\sigma}{\sqrt n},\ \Xbar+z_{1-\alpha/2}\frac{\sigma}{\sqrt n}\Bigr]
\]
tiene cobertura exactamente $1-\alpha$. Si los datos no son normales pero $n$ es grande, la cobertura es aproximadamente $1-\alpha$.
\end{teorema}

La cantidad $E=z_{1-\alpha/2}\,\sigma/\sqrt n$ es el \emph{margen de error} (la semiamplitud del intervalo). Los valores más usados son $z_{0{,}95}=1{,}645$ (nivel $90\,\%$), $z_{0{,}975}=1{,}960$ ($95\,\%$) y $z_{0{,}995}=2{,}576$ ($99\,\%$).

\begin{corolario}[Tamaño muestral]\label{cor:t4-tamano}
Para que el margen de error del \cref{teo:t4-ic-z} sea a lo sumo $E$ basta tomar
\[
n\ge\Bigl(\frac{z_{1-\alpha/2}\,\sigma}{E}\Bigr)^{2},
\]
redondeando al entero superior.
\end{corolario}

\section{Intervalo para la media con \texorpdfstring{$\sigma$}{sigma} desconocida}

En la práctica $\sigma$ casi nunca se conoce. Al sustituirla por $S$ hay que reemplazar la normal por la $t$ de Student.

\begin{teorema}[Intervalo $t$]\label{teo:t4-ic-t}
Sean $X_1,\dots,X_n\sim\Nor(\mu,\sigma^{2})$ \iid\ con $\sigma$ desconocida y $n\ge2$. El intervalo
\[
\Bigl[\Xbar-t_{n-1;1-\alpha/2}\frac{S}{\sqrt n},\ \Xbar+t_{n-1;1-\alpha/2}\frac{S}{\sqrt n}\Bigr]
\]
tiene cobertura exactamente $1-\alpha$.
\end{teorema}

Como $t_{n-1;1-\alpha/2}>z_{1-\alpha/2}$, el intervalo $t$ es más ancho: refleja la incertidumbre adicional por estimar $\sigma$. Para $n$ grande (digamos $n\ge30$ o más) ambos cuantiles casi coinciden y, por el teorema central del límite, el intervalo es aproximadamente válido aunque los datos no sean normales.

\addtocounter{definicion}{1}%

\section{Intervalos para proporciones}

Si $X_1,\dots,X_n\sim\Ber(p)$, el \cref{cor:t3-binomial} da $\widehat p\approx\Nor\bigl(p,p(1-p)/n\bigr)$. El error estándar depende de $p$, desconocido, y se estima por $\sqrt{\widehat p(1-\widehat p)/n}$.

\begin{definicion}[Intervalo de Wald para una proporción]\label{def:t4-wald}
El intervalo de confianza aproximado de nivel $1-\alpha$ para $p$ es
\[
\widehat p\pm z_{1-\alpha/2}\sqrt{\frac{\widehat p(1-\widehat p)}{n}}.
\]
Como $p(1-p)\le1/4$, el intervalo \emph{conservador} $\widehat p\pm z_{1-\alpha/2}/(2\sqrt n)$ es siempre al menos tan ancho~\cite[§8.3]{openstax}.
\end{definicion}

La aproximación es buena si $n\widehat p$ y $n(1-\widehat p)$ son ambos al menos del orden de $10$; con muestras pequeñas o con $\widehat p$ cercano a $0$ o a $1$ la cobertura real es inferior a la nominal. Existen correcciones sencillas, como el método de «más cuatro» de~\cite[§8.3]{openstax} (sumar dos éxitos y dos fracasos antes de calcular) o el intervalo de Wilson, que quedan fuera de este módulo.

\begin{corolario}[Tamaño muestral para una proporción]\label{cor:t4-tamano-p}
Para que el margen de error sea a lo sumo $E$ sea cual sea $p$, basta $n\ge z_{1-\alpha/2}^{2}/(4E^{2})$. Si se dispone de una estimación previa $p^{*}$, basta $n\ge z_{1-\alpha/2}^{2}\,p^{*}(1-p^{*})/E^{2}$.
\end{corolario}

\addtocounter{definicion}{1}%

\addtocounter{definicion}{1}%

\addtocounter{definicion}{1}%


% =====================================================================
\chapter{Contrastes de hipótesis}
\label{tema:5}
% =====================================================================

Un contraste de hipótesis decide, a partir de los datos, si hay evidencia suficiente contra una afirmación sobre la población: que una moneda es equilibrada, que un nuevo modelo no mejora al anterior o que unos recuentos siguen una distribución dada. Este tema presenta el marco de Neyman y Pearson (hipótesis nula y alternativa, región de rechazo, errores de tipo~I y~II, potencia), el $p$-valor, los contrastes $z$ y $t$ para medias y proporciones, el contraste $\chi^{2}$ de bondad de ajuste y el problema de las comparaciones múltiples, con su aplicación a la evaluación de modelos de lenguaje en muchos benchmarks. Se sigue el capítulo~10 de Wasserman~\cite{wasserman}, «Hypothesis Testing and p-values», las lecturas~17b, 18 y~19 del curso 18.05 del MIT~\cite{mit1805}, las clases~5, 7 y~11 de Stats~200~\cite{stats200} y los capítulos~9 y~11 de OpenStax~\cite{openstax}.

\section{Planteamiento de un contraste}

\begin{definicion}[Hipótesis, estadístico y región de rechazo]\label{def:t5-contraste}
Sea $\{F_\theta:\theta\in\Theta\}$ un modelo y $\Theta_0\subset\Theta$. Se contrastan la \emph{hipótesis nula} $H_0:\theta\in\Theta_0$ frente a la \emph{hipótesis alternativa} $H_1:\theta\in\Theta\setminus\Theta_0$. Una hipótesis es \emph{simple} si especifica un único valor de $\theta$ y \emph{compuesta} en otro caso. Un \emph{contraste} consiste en un \emph{estadístico de contraste} $T=T(X_1,\dots,X_n)$ y una \emph{región de rechazo} $R$: si $T\in R$ se rechaza $H_0$; si no, no se rechaza.
\end{definicion}

La hipótesis nula representa la situación por defecto («no hay efecto», «la moneda es equilibrada», «los dos modelos son igual de buenos») y el contraste busca evidencia contra ella. Las alternativas más comunes para un parámetro real son la \emph{bilateral} $H_1:\theta\neq\theta_0$ y las \emph{unilaterales} $H_1:\theta>\theta_0$ o $H_1:\theta<\theta_0$.

\begin{definicion}[Errores, nivel y potencia]\label{def:t5-errores}
Se comete un \emph{error de tipo~I} si se rechaza $H_0$ siendo cierta, y un \emph{error de tipo~II} si no se rechaza siendo falsa:
\begin{center}
\small
\begin{tabular}{@{}lcc@{}}
\toprule
 & $H_0$ cierta & $H_0$ falsa \\
\midrule
No se rechaza $H_0$ & decisión correcta & error de tipo~II \\
Se rechaza $H_0$ & error de tipo~I & decisión correcta \\
\bottomrule
\end{tabular}
\end{center}
El contraste tiene \emph{nivel de significación} $\alpha$ si $\PP_\theta(T\in R)\le\alpha$ para todo $\theta\in\Theta_0$. La \emph{función de potencia} es $\pi(\theta)=\PP_\theta(T\in R)$; para $\theta\in\Theta\setminus\Theta_0$, $\pi(\theta)$ es la \emph{potencia} contra esa alternativa y $\beta(\theta)=1-\pi(\theta)$ es la probabilidad de error de tipo~II~\cite[cap.~10]{wasserman}.
\end{definicion}

Los dos errores no se pueden reducir a la vez con un tamaño muestral fijo: agrandar la región de rechazo reduce los errores de tipo~II y aumenta los de tipo~I. El convenio es fijar $\alpha$ (habitualmente $0{,}05$ o $0{,}01$) y, entre los contrastes de nivel $\alpha$, preferir los más potentes. Las dos hipótesis no son simétricas: no rechazar $H_0$ no prueba que sea cierta, solo que los datos no aportan evidencia suficiente contra ella.

\section{El \texorpdfstring{$p$}{p}-valor}

\begin{definicion}[$p$-valor]\label{def:t5-pvalor}
Supóngase que se rechaza $H_0$ para valores grandes de un estadístico $T$ y que $H_0$ es simple. Si el valor observado es $t_{\mathrm{obs}}$, el \emph{$p$-valor} es
\[
p=\PP_{H_0}(T\ge t_{\mathrm{obs}}),
\]
la probabilidad, calculada suponiendo $H_0$ cierta, de obtener un resultado al menos tan extremo como el observado. (Si se rechaza para valores pequeños, $p=\PP_{H_0}(T\le t_{\mathrm{obs}})$; en un contraste bilateral con estadístico de distribución simétrica, $p=\PP_{H_0}(\abs T\ge\abs{t_{\mathrm{obs}}})$.)
\end{definicion}

\begin{proposicion}[Decisión con el $p$-valor]\label{prop:t5-pvalor}
Sea un contraste que rechaza $H_0$ si $T\ge c_\alpha$, donde $c_\alpha$ cumple $\PP_{H_0}(T\ge c_\alpha)=\alpha$ y $T$ tiene distribución continua bajo $H_0$. Entonces se rechaza $H_0$ al nivel $\alpha$ si y solo si $p\le\alpha$.
\end{proposicion}

El $p$-valor es, por tanto, el menor nivel al que los datos permiten rechazar $H_0$, y resume en un número la fuerza de la evidencia contra $H_0$: cuanto menor, más incompatibles son los datos con $H_0$.

\begin{observacion}[Lo que el $p$-valor no es]\label{obs:t5-pvalor}
El $p$-valor \emph{no} es la probabilidad de que $H_0$ sea cierta, ni la probabilidad de que los resultados se deban al azar, ni mide el tamaño del efecto. Un $p$-valor muy pequeño puede acompañar a un efecto irrelevante si $n$ es enorme, y un efecto grande puede dar un $p$-valor alto si $n$ es pequeño. Por eso conviene informar siempre del efecto estimado con su intervalo de confianza, además del $p$-valor~\cite[lecturas~17b y~20]{mit1805}.
\end{observacion}

\section{Contrastes \texorpdfstring{$z$}{z} y \texorpdfstring{$t$}{t}}

Los contrastes de esta sección usan como estadístico una estimación tipificada con su error estándar calculado bajo $H_0$. En todos ellos, las regiones de rechazo de nivel $\alpha$ son $T>q_{1-\alpha}$ (alternativa $>$), $T<-q_{1-\alpha}$ (alternativa $<$) y $\abs T>q_{1-\alpha/2}$ (bilateral), donde $q$ es el cuantil de la distribución de referencia.

\begin{center}
\small
\begin{tabular}{@{}llll@{}}
\toprule
Contraste & $H_0$ & Estadístico & Distribución bajo $H_0$ \\
\midrule
$z$ para la media ($\sigma$ conocida) & $\mu=\mu_0$ & $Z=\dfrac{\Xbar-\mu_0}{\sigma/\sqrt n}$ & $\Nor(0,1)$ (datos normales) \\[2ex]
$t$ para la media & $\mu=\mu_0$ & $T=\dfrac{\Xbar-\mu_0}{S/\sqrt n}$ & $t_{n-1}$ (datos normales) \\[2ex]
$z$ para una proporción & $p=p_0$ & $Z=\dfrac{\widehat p-p_0}{\sqrt{p_0(1-p_0)/n}}$ & $\approx\Nor(0,1)$ ($n$ grande) \\
\bottomrule
\end{tabular}
\end{center}

La distribución bajo $H_0$ se sigue del \cref{teo:t3-suma-normales}, del \cref{cor:t3-t} y del \cref{cor:t3-binomial}, respectivamente. En el contraste de una proporción, el error estándar se calcula con $p_0$, no con $\widehat p$, porque se razona bajo $H_0$.

\addtocounter{definicion}{1}%

\section{Potencia y tamaño muestral}

\begin{proposicion}[Potencia del contraste $z$ unilateral]\label{prop:t5-potencia}
Para el contraste $z$ de $H_0:\mu=\mu_0$ frente a $H_1:\mu>\mu_0$ al nivel $\alpha$ con datos $\Nor(\mu,\sigma^{2})$, la función de potencia es
\[
\pi(\mu)=1-\Phi\Bigl(z_{1-\alpha}-\frac{(\mu-\mu_0)\sqrt n}{\sigma}\Bigr).
\]
\end{proposicion}

En el \cref{ej:t5-z}, si la media real fuera $220$, la potencia sería $1-\Phi(1{,}645-2)=\Phi(0{,}355)\approx0{,}64$: con $64$ respuestas, el contraste detectaría esa diferencia solo en dos de cada tres ocasiones. El \cref{eje:t5-potencia} deduce el tamaño muestral necesario para alcanzar una potencia dada.

\section{Contraste \texorpdfstring{$\chi^{2}$}{chi cuadrado} de bondad de ajuste}

Sea una variable categórica con $k$ categorías y probabilidades $p_1,\dots,p_k$. En $n$ observaciones independientes se obtienen los recuentos $O_1,\dots,O_k$, con $\sum_jO_j=n$. Se quiere contrastar $H_0:p_j=p_{0j}$ para todo $j$, con $p_{0j}$ dados.

\begin{definicion}[Estadístico $\chi^{2}$ de Pearson]\label{def:t5-chi2}
Sean $E_j=np_{0j}$ los recuentos \emph{esperados} bajo $H_0$. El estadístico es
\[
X^{2}=\sum_{j=1}^{k}\frac{(O_j-E_j)^{2}}{E_j}.
\]
Bajo $H_0$, y si todos los $E_j$ son al menos~$5$, $X^{2}$ sigue aproximadamente una distribución $\chi^{2}_{k-1}$. Se rechaza $H_0$ al nivel $\alpha$ si $X^{2}>\chi^{2}_{k-1;1-\alpha}$~\cite[§11.2]{openstax}.
\end{definicion}

Los $k-1$ grados de libertad se deben a la restricción $\sum_jO_j=n$: conocidos $k-1$ recuentos, el último queda determinado. Si se estiman $m$ parámetros del modelo a partir de los mismos datos, se restan también: $k-1-m$ grados de libertad. Si algún $E_j$ es menor que $5$, se agrupan categorías contiguas.

\addtocounter{definicion}{1}%

\section{Comparaciones múltiples}

Si se realizan muchos contrastes, cada uno al nivel $\alpha$, la probabilidad de cometer \emph{algún} error de tipo~I crece rápidamente.

\begin{proposicion}\label{prop:t5-multiples}
Si se realizan $m$ contrastes independientes al nivel $\alpha$ y todas las hipótesis nulas son ciertas, la probabilidad de rechazar al menos una es $1-(1-\alpha)^{m}$, y el número esperado de rechazos (todos falsos) es $m\alpha$.
\end{proposicion}

Con $m=20$ y $\alpha=0{,}05$, la probabilidad de al menos un «descubrimiento» falso es $1-0{,}95^{20}\approx0{,}64$.

\begin{definicion}[Tasa de error por familia]\label{def:t5-fwer}
La \emph{tasa de error por familia} (FWER, \emph{family-wise error rate}) de un procedimiento que realiza $m$ contrastes es la probabilidad de cometer al menos un error de tipo~I~\cite[clase~11]{stats200}.
\end{definicion}

\begin{teorema}[Corrección de Bonferroni]\label{teo:t5-bonferroni}
Si cada uno de los $m$ contrastes se realiza al nivel $\alpha/m$ (es decir, se rechaza la hipótesis $i$ si $p_i\le\alpha/m$), la tasa de error por familia es a lo sumo $\alpha$, cualquiera que sea la dependencia entre los contrastes y cualquiera que sea el conjunto de hipótesis nulas ciertas.
\end{teorema}

La corrección de Bonferroni es sencilla y muy general, pero conservadora: con $m$ grande reduce mucho la potencia. Existen procedimientos más potentes, como el de Holm o el de Benjamini y Hochberg, que controla una cantidad distinta, la tasa de falsos descubrimientos~\cite[clase~11]{stats200}; quedan fuera de este módulo.

\begin{observacion}[Comparar modelos en muchos benchmarks]\label{obs:t5-benchmarks}
La evaluación de modelos de lenguaje es un terreno abonado para las comparaciones múltiples: se comparan dos modelos en decenas de benchmarks, se prueban muchas variantes de \emph{prompt} y se informa de la mejor, o se ordena una tabla con cientos de modelos. Si en $20$ benchmarks independientes no hay ninguna diferencia real, se esperan $20\cdot0{,}05=1$ diferencias «significativas» al $5\,\%$. Para comparar la exactitud de dos modelos se puede usar un contraste $z$ de diferencia de proporciones (\cref{eje:t5-dos-modelos}); cuando ambos modelos responden a las \emph{mismas} preguntas, es preferible un contraste pareado que use las diferencias pregunta a pregunta, porque elimina la variabilidad debida a la dificultad de cada pregunta (ampliación).
\end{observacion}

\addtocounter{definicion}{1}%

\addtocounter{definicion}{1}%


% =====================================================================
\chapter{Regresión lineal simple y correlación}
\label{tema:6}
% =====================================================================

La regresión estudia cómo cambia en promedio una variable respuesta $y$ cuando cambia una variable explicativa $x$. Este tema deduce la recta de mínimos cuadrados, estudia los residuos y la descomposición de la variabilidad que da lugar al coeficiente de determinación $R^{2}$, muestra cómo las transformaciones logarítmicas permiten ajustar leyes de potencia y termina con la distinción entre correlación y causalidad. Se sigue el capítulo~12 de OpenStax~\cite{openstax} (§12.3 «The Regression Equation»), el capítulo~14 de Rice~\cite{rice}, «Linear Least Squares», el capítulo~13 de Wasserman~\cite{wasserman}, «Linear and Logistic Regression», y la lectura~26 del curso 18.05 del MIT~\cite{mit1805}; para la causalidad, el capítulo~16 de Wasserman, «Causal Inference».

\section{La recta de mínimos cuadrados}

Se dispone de $n\ge2$ pares $(x_1,y_1),\dots,(x_n,y_n)$ con los $x_i$ no todos iguales. Se usan las sumas centradas
\[
S_{xx}=\sum_{i=1}^{n}(x_i-\xbar)^{2},\qquad S_{yy}=\sum_{i=1}^{n}(y_i-\ybar)^{2},\qquad S_{xy}=\sum_{i=1}^{n}(x_i-\xbar)(y_i-\ybar),
\]
de modo que $s_x^{2}=S_{xx}/(n-1)$, $s_{xy}=S_{xy}/(n-1)$ y $r=S_{xy}/\sqrt{S_{xx}S_{yy}}$ (\cref{def:t1-cov}). Por la \cref{prop:t1-abreviada}, $S_{xx}=\sum x_i^{2}-n\xbar^{2}$ y, análogamente, $S_{xy}=\sum x_iy_i-n\xbar\ybar$.

\begin{definicion}[Recta de mínimos cuadrados]\label{def:t6-mc}
Para cada recta $y=b_0+b_1x$, el \emph{valor ajustado} del dato $i$ es $\widehat y_i=b_0+b_1x_i$ y el \emph{residuo} es $e_i=y_i-\widehat y_i$. La \emph{recta de mínimos cuadrados} (o de regresión de $y$ sobre $x$) es la que minimiza la suma de cuadrados de los residuos
\[
\operatorname{SC}(b_0,b_1)=\sum_{i=1}^{n}(y_i-b_0-b_1x_i)^{2}.
\]
\end{definicion}

\begin{teorema}[Coeficientes de mínimos cuadrados]\label{teo:t6-mc}
Si $S_{xx}>0$, la función $\operatorname{SC}$ tiene un único mínimo, en
\[
\widehat b_1=\frac{S_{xy}}{S_{xx}}=r\,\frac{s_y}{s_x},\qquad\widehat b_0=\ybar-\widehat b_1\xbar,
\]
y el valor mínimo es $S_{yy}-S_{xy}^{2}/S_{xx}=S_{yy}(1-r^{2})$.
\end{teorema}

La condición $\widehat b_0=\ybar-\widehat b_1\xbar$ dice que la recta pasa por el punto de medias $(\xbar,\ybar)$. La pendiente tiene el signo de $r$: es positiva cuando la asociación es creciente.

\addtocounter{definicion}{1}%

\addtocounter{figure}{1}%

\section{Residuos y coeficiente de determinación}

\begin{proposicion}[Propiedades de los residuos]\label{prop:t6-residuos}
Para la recta de mínimos cuadrados, $\sum_ie_i=0$, $\sum_ix_ie_i=0$ y $\sum_i\widehat y_ie_i=0$.
\end{proposicion}

La demostración es el \cref{eje:t6-residuos}. Las dos primeras igualdades son las \emph{ecuaciones normales} (\cref{eje:t6-normales}): los residuos son ortogonales, como vectores de $\R^{n}$, al vector de unos y al vector de las $x_i$. Este es el punto de vista geométrico que el módulo~8 generaliza a varias variables explicativas.

\begin{teorema}[Descomposición de la variabilidad]\label{teo:t6-anova}
Sean $\operatorname{SC}_T=\sum_i(y_i-\ybar)^{2}=S_{yy}$ (suma de cuadrados total), $\operatorname{SC}_R=\sum_i(\widehat y_i-\ybar)^{2}$ (explicada por la regresión) y $\operatorname{SC}_E=\sum_ie_i^{2}$ (residual o de error). Entonces
\[
\operatorname{SC}_T=\operatorname{SC}_R+\operatorname{SC}_E.
\]
\end{teorema}

\begin{definicion}[Coeficiente de determinación]\label{def:t6-r2}
El \emph{coeficiente de determinación} es $R^{2}=\operatorname{SC}_R/\operatorname{SC}_T=1-\operatorname{SC}_E/\operatorname{SC}_T$, la fracción de la variabilidad de $y$ explicada por la recta. Por el \cref{teo:t6-anova}, $0\le R^{2}\le1$.
\end{definicion}

\begin{corolario}\label{cor:t6-r2}
En la regresión lineal simple, $R^{2}=r^{2}$.
\end{corolario}

En el \cref{ej:t6-recta}, $R^{2}=0{,}6$: la recta explica el $60\,\%$ de la variabilidad de $y$. La dispersión de los puntos alrededor de la recta se resume en el \emph{error típico residual} $s_e=\sqrt{\operatorname{SC}_E/(n-2)}$; se divide por $n-2$ porque se han estimado dos parámetros (compárese con la \cref{obs:t1-n-1}).

\section{Diagnóstico con residuos}

Un $R^{2}$ alto no garantiza que el modelo lineal sea adecuado. La herramienta básica de diagnóstico es el \emph{gráfico de residuos}, que representa $e_i$ frente a $x_i$ (o frente a $\widehat y_i$). Si el modelo es adecuado, los residuos se distribuyen sin estructura alrededor de cero. Algunas señales de alarma~\cite[§12.3]{openstax}:
\begin{enumerate}
  \item una curvatura sistemática (residuos positivos en los extremos y negativos en el centro, o al revés) indica que la relación no es lineal;
  \item una dispersión que crece con $x$ (forma de embudo) indica varianza no constante;
  \item un punto con un residuo muy grande es un posible atípico, y un punto con $x_i$ muy alejado de $\xbar$ puede ser muy \emph{influyente}: puede cambiar mucho la recta por sí solo.
\end{enumerate}
Además, la recta solo describe la relación en el rango de los $x_i$ observados: \emph{extrapolar} fuera de ese rango (\cref{eje:t6-buceo}) es arriesgado.

\begin{observacion}[El modelo estadístico]\label{obs:t6-modelo}
Si se supone que $y_i=\beta_0+\beta_1x_i+\varepsilon_i$ con errores $\varepsilon_i\sim\Nor(0,\sigma^{2})$ \iid, la log-verosimilitud de $(\beta_0,\beta_1)$ es una constante menos $\frac{1}{2\sigma^{2}}\operatorname{SC}(\beta_0,\beta_1)$; por tanto, los estimadores de mínimos cuadrados son los estimadores de máxima verosimilitud~\cite[cap.~13]{wasserman}, \cite[\emph{problem set}~6]{mit1805}. Con este modelo pueden construirse intervalos de confianza y contrastes para la pendiente (ampliación), que se estudian en~\cite[cap.~14]{rice}.
\end{observacion}

\section{Transformaciones: ajuste de leyes de potencia}

Muchas relaciones no lineales se vuelven lineales tras una transformación. Una \emph{ley de potencia} $y=a\,x^{b}$ con $a,x,y>0$ cumple
\[
\log y=\log a+b\log x,
\]
de modo que en escala log-log es una recta de pendiente $b$ y ordenada $\log a$. Para ajustarla se aplica mínimos cuadrados a los pares $(\log x_i,\log y_i)$ y se deshace la transformación: $\widehat a=10^{\widehat b_0}$ (o $e^{\widehat b_0}$ si se usan logaritmos naturales) y $\widehat b=\widehat b_1$. Análogamente, una relación exponencial $y=a\,e^{bx}$ se linealiza con $\ln y=\ln a+bx$. Obsérvese que así se minimizan los errores \emph{relativos} (en escala logarítmica), no los absolutos.

\addtocounter{definicion}{1}%

\section{Correlación y causalidad}

Una correlación entre $x$ e $y$ puede deberse a que $x$ influye en $y$, a que $y$ influye en $x$ (causalidad inversa), a que una tercera variable $z$ influye en ambas (\emph{variable de confusión}) o simplemente al azar. Los datos observacionales, por sí solos, no permiten distinguir estos casos~\cite[cap.~16]{wasserman}.

\addtocounter{definicion}{1}%

Un caso extremo de confusión es la \emph{paradoja de Simpson}: una asociación que se da en cada uno de los grupos de una población puede invertirse al agregar los grupos, cuando los tamaños de los grupos difieren entre las condiciones comparadas (\cref{eje:t6-simpson}).

La forma más fiable de establecer una relación causal es el \emph{experimento aleatorizado}: si el tratamiento (por ejemplo, el \emph{prompt} A o el B) se asigna al azar a las unidades (las preguntas), las variables de confusión se reparten en promedio por igual entre los grupos, y una diferencia sistemática en la respuesta solo puede atribuirse al tratamiento o al azar, cuya magnitud se cuantifica con los contrastes del \cref{tema:5}.

\addtocounter{definicion}{1}%

\addtocounter{definicion}{1}%



\backmatter
% =====================================================================
\begin{thebibliography}{99}
\addcontentsline{toc}{chapter}{Bibliografía}
\markboth{Bibliografía}{}
\small

\bibitem{bh}
J.~K. Blitzstein y J.~Hwang.
\newblock \emph{Introduction to Probability}, 2.ª ed.
\newblock CRC Press, 2019.
\newblock Texto del curso Stat~110 de Harvard. Capítulos~4 «Expectation», 5 «Continuous Random Variables», 6 «Moments», 7 «Joint Distributions» y 10 «Inequalities and Limit Theorems».

\bibitem{openstax}
B.~Illowsky y S.~Dean.
\newblock \emph{Introductory Statistics}, 2.ª ed.
\newblock OpenStax, Rice University, 2023. Libro de texto de acceso libre. Secciones 2.2--2.7, 7.1--7.2, 8.1--8.3, 9.1--9.5, 11.2 y 12.3--12.4.
\newblock \url{https://openstax.org/details/books/introductory-statistics-2e}

\bibitem{mit1805}
J.~Orloff y J.~Bloom.
\newblock 18.05: Introduction to Probability and Statistics.
\newblock Massachusetts Institute of Technology, MIT OpenCourseWare, primavera de 2022. Lecturas (\emph{readings}) 5c, 6b, 7b, 10a, 10b, 17b, 18, 19, 20, 22, 23a, 23b y 26; \emph{problem sets}~4, 5, 6, 8, 9, 10 y 11; problemas de la clase~23.
\newblock \url{https://ocw.mit.edu/courses/18-05-introduction-to-probability-and-statistics-spring-2022/}

\bibitem{rice}
J.~A. Rice.
\newblock \emph{Mathematical Statistics and Data Analysis}, 3.ª ed.
\newblock Duxbury (Thomson Brooks/Cole), 2007. Capítulos~8 «Estimation of Parameters and Fitting of Probability Distributions», 9 «Testing Hypotheses and Assessing Goodness of Fit», 10 «Summarizing Data» y 14 «Linear Least Squares».

\bibitem{stats200}
Z.~Fan.
\newblock Stats~200: Introduction to Statistical Inference.
\newblock Stanford University, otoño de 2016 (edición archivada \texttt{stats200.1172}). Clases~5, 7, 11, 12, 13, 18 y 25; \emph{homework}~3, 4 y~5.
\newblock \url{https://web.stanford.edu/class/archive/stats/stats200/stats200.1172/}

\bibitem{wasserman}
L.~Wasserman.
\newblock \emph{All of Statistics: A Concise Course in Statistical Inference}.
\newblock Springer, 2004. Capítulos~5 «Convergence of Random Variables», 6 «Models, Statistical Inference and Learning», 9 «Parametric Inference», 10 «Hypothesis Testing and p-values», 13 «Linear and Logistic Regression» y 16 «Causal Inference».

\end{thebibliography}


\end{document}
